% concatenated from V5_Commutant_lifting_and_interpolation_on_quotients_of_bounded_symmetric_domains.tex
 %\documentclass[draft,reqno, 10pt]{amsproc}
%\documentclass[a4paper,oneside,11pt, reqno]{amsproc}
\documentclass[12pt,reqno]{amsart}
\usepackage{amsthm}
\usepackage{amsfonts}
\usepackage{amsmath}
\usepackage{amssymb}
\usepackage{mathrsfs}
\usepackage{epsfig}
\usepackage{graphicx}
\usepackage{epstopdf}
%\usepackage{pst-node}
\usepackage{tikz-cd}
\usepackage{color}
\usepackage{ifthen}
\usepackage{float}
\usepackage{cancel} 
\usepackage{comment}
\usepackage{stmaryrd}
\usepackage{hyperref}
%\usepackage{refcheck}
\makeatletter
\@namedef{subjclassname@2010}{%
\textup{2010} Mathematics Subject Classification}
\makeatother
\usepackage[normalem]{ulem}
\newtheorem{theorem}{Theorem}[section]
\newtheorem{proposition}[theorem]{Proposition}%[section]
\newtheorem{corollary}[theorem]{Corollary}%[section]
\newtheorem{lemma}[theorem]{Lemma}%[section]
\newtheorem*{Convention}{Convention}
\newtheorem{question}[theorem]{\bf Question}
\newtheorem{remark}[theorem]{\bf Remark}%[section]

\setlength{\textheight}{680pt} \setlength{\textwidth}{450pt}
\oddsidemargin -0mm \evensidemargin -0mm \topmargin -25pt
\baselineskip=0.30in

\theoremstyle{remark}

% \newtheorem{dummy}{Theorem} % dummy counter
% \newenvironment{realizationthm}[1][]
%   {\renewcommand\thedummy{R}%
%    \begin{dummy}[\bf Realization theorem]}
%   {\end{dummy}}

%  \newtheorem{dummyy}{Theorem} % dummy counter
% \newenvironment{interpolationthm}[1][]
%   {\renewcommand\thedummyy{I}%
%    \begin{dummyy}[\bf Interpolation theorem]}
%   {\end{dummyy}}
  
%   \newtheorem{dummyyy}{Theorem} % dummy counter
% \newenvironment{Toeplitzcoronathm}[1][]
%   {\renewcommand\thedummyyy{T-C}%
%    \begin{dummyyy}[\bf Toeplitz corona theorem]}
%   {\end{dummyyy}}
  
%   \newtheorem{dummyyyy}{Theorem} % dummy counter
% \newenvironment{extensionthm}[1][]
%   {\renewcommand\thedummyyyy{E}%
%    \begin{dummyyyy}[\bf Extension theorem]}
%   {\end{dummyyyy}}
 \newcommand{\cred}{\color{red}}
\newcommand{\cpl}{\color{purple}}
\newcommand{\cblue}{\color{blue}}
\theoremstyle{definition}
\newtheorem{definition}[theorem]{Definition}%[section]
\newtheorem{example}[theorem]{Example}
\newtheorem{List}[theorem]{List}
\newtheorem*{Caution}{Caution}
\def \dep{\mathsf{d}}
\newcommand*\initfamily{\usefont{U}{Carrickc}{xl}{n}}
\newcommand{\ran}{\mbox{ran}}
\newcommand{\ol}{\overline}
\newcommand{\bq}{\begin{equation}}
\newcommand{\eq}{\end{equation}}
\newcommand{\beqn}{\begin{eqnarray*}}
\newcommand{\eeqn}{\end{eqnarray*}}
\newcommand{\beq}{\begin{eqnarray}}
\newcommand{\eeq}{\end{eqnarray}}
\newcommand{\nno}{\nonumber}
\usepackage{soul}
\newcommand{\rar}{\rightarrow}
\newcommand{\Rar}{\Rightarrow}
\newcommand{\noin}{\noindent}
\newcommand{\bc}{\begin{centre}}
\newcommand{\ec}{\end{centre}}
\newcommand{\sz}{\scriptsize}
\newcommand{\rf}{\ref}
\newcommand{\sgm}{\sigma}
\newcommand{\Sgm}{\Sigma}
\newcommand{\dt}{\delta}
% \newcommand{\SAE}{\mathcal S \mathcal A(\mathbb E)} \newcommand{\AKE}{\mathcal A\mathcal K (\mathbb E)}
% \newcommand{\UCE}{\mathcal U \mathcal C(\mathbb E)}
% \newcommand{\CDF}{C(\overline{\mathbb D})_F^+}
% \newcommand{\CDE}{C(\overline{\mathbb D})_{\mathbb E}^+}
% \newcommand{\PKF}{\mathbb {C}_F^+}
% \newcommand{\PKE}{\mathbb {C}_{\mathbb E}^+}
% \newcommand{\SAFE}{\mathcal{SA}_f(\mathbb E)}
\newcommand{\Dt}{Delta}
\newcommand{\lmd}{\lambda}
\newcommand{\Lmd}{\Lambda}
\newcommand{\eps}{\epsilon}
\newcommand{\pcc}{\stackrel{P}{>}}
\newcommand{\dist}{{\rm\,dist}}
\newcommand{\sspan}{{\rm\,span}}
\newcommand{\re}{{\rm Re\,}}
\newcommand{\im}{{\rm Im\,}}
\newcommand{\sgn}{{\rm sgn\,}}
\newcommand{\ba}{\begin{array}}
\newcommand{\ea}{\end{array}}
\newcommand{\eop}{\hfill{{\rule{2.5mm}{2.5mm}}}}
\newcommand{\eoe}{\hfill{{\rule{1.5mm}{1.5mm}}}}
\newcommand{\eof}{\hfill{{\rule{1.5mm}{1.5mm}}}}
\newcommand{\hone}{\mbox{\hspace{1em}}}
\newcommand{\htwo}{\mbox{\hspace{2em}}}
\newcommand{\hthree}{\mbox{\hspace{3em}}}
\newcommand{\hfour}{\mbox{\hspace{4em}}}
\newcommand{\hsev}{\mbox{\hspace{7em}}}
\newcommand{\vone}{\vskip 2ex}
\newcommand{\vtwo}{\vskip 4ex}
\newcommand{\vonee}{\vskip 1.5ex}
\newcommand{\vthree}{\vskip 6ex}
\newcommand{\vfour}{\vspace*{8ex}}
\newcommand{\norm}{\|\;\;\|}
\newcommand{\integ}[4]{\int_{#1}^{#2}\,{#3}\,d{#4}}
\newcommand{\inp}[2]{\langle{#1},\,{#2} \rangle}
\newcommand{\Inp}[2]{\Langle{#1},\,{#2} \Langle}
\newcommand{\vspan}[1]{{{\rm\,span}\#1 \}}}
\newcommand{\dm}[1]{\displaystyle {#1}}
\newcommand*{\childi}[2]{{\mathsf{Chi}}_{#1}(#2)}
\newcommand*{\parenti}[2]{{\mathsf{par}}_{#1}(#2)}
\newcommand{\Shubhamco}[1]{{\color{blue}[Shubham:~#1]}}
\newcommand{\Shubhamcom}[1]{{\color{purple}[Shubham:~#1]}}
\newcommand{\Paramitaco}[1]{{\color{purple}[[Paramita:~#1]]}}
\newcommand{\Milanco}[1]{{\color{red}[[Milan:~#1]]}}
\newcommand{\Surjitco}[1]{{\color{brown}[[Surjit:~#1]]}}
\renewcommand{\Delta}{{\nabla}}
\def \N{\mathbb{Z}_+}
\def \D{\mathbb{D}}
\def \R{\mathbb{R}}
\def \C{\mathbb{C}}
\def \F{\mathbb{F}}
\def \T{\mathbb{T}}
\def \s{\boldsymbol{s}}
\def \z{\boldsymbol{z}}
\def \\theta{\boldsymbol{\theta}}
\def \\xi{\boldsymbol{\xi}}
\def \q{\boldsymbol{q}}
\def \anti{\mathrm{anti}}
\def \sym{\mathrm{sym}}

\newcommand{\mf}{\mathfrak}
\newcommand*{\des}[1]{{{\mathsf{Des}}(#1)}}
\newcommand*{\dess}[2]{{{\mathsf{Des}}_{#1}(#2)}}
\newcommand*{\child}[1]{\mathsf{Chi}(#1)}
\newcommand*{\childn}[2]{{\mathsf{Chi}}^{\langle#1\rangle}(#2)}
\newcommand*{\parentn}[2]{{\mathsf{par}}^{\langle#1\rangle}(#2)}
\newcommand*{\gammab}{\boldsymbol\gamma}
\newcommand*{\Ge}{\geqslant}
\newcommand*{\lambdab}{\boldsymbol\lambda}
\newcommand*{\Le}{\leqslant}
\newcommand*{\parent}[1]{\mathsf{par}(#1)}
\newcommand*{\rot}{\mathsf{\omega}}
%\newcommand*{S_{\lambda}}{S_{\lambdab}}
\newcommand*{\spr}[1]{\color{green}{#1}\color{black}\/}
\newcommand*{\suc}[1]{\mathsf{Suc}(#1)}
\newcommand*{\supp}{\mathrm{supp}}
%\newcommand*{{\mathscr T}}{\mathscr T}
\newcommand*{\uwam}[1]{\marginpar{\sf\tiny{{\textcolor{red}
{$\blacktriangleleft$} \textcolor{blue}{#1}
\textcolor{red}{$\blacktriangleright$}}}\normalsize}}


\begin{document}
\title[Commutant lifting]{Commutant lifting and interpolation on quotients of bounded symmetric domains}

\author[M. K. Mal]{Milan Kumar Mal}
\address[M. K. Mal]{Department of Mathematics, Indian Institute of Technology Madras, Chennai 600036, India}
 \email{ma21d018@smail.iitm.ac.in; milanmal1702@gmail.com }


\thanks{Support for the work of the author was provided in the form of a Prime Minister's Research Fellowship (PMRF / 2502827).}

 \subjclass[2020]{Primary 47A57, 47A20, 47B32 Secondary 32H35, 32E30}

 \keywords{Commutant lifting, Nevanlinna-Pick interpolation, quotient domain, bounded symmetric domain, complex reflection group, Hardy space, inner function}


%\thanks{}


\date{}

\begin{abstract}
Let $\Omega\subseteq \mathbb C^d$ be a bounded symmetric domain, $G$ a finite complex reflection group acting on $\mathbb C^d$, and $\boldsymbol \theta:\Omega\to \boldsymbol \theta(\Omega)$ the associated proper holomorphic map factored by $G.$ In this paper, we investigate commutant lifting and interpolation by Schur functions on the quotient domain $\boldsymbol \theta(\Omega).$
For a given quotient module of the Hardy space $H^2(\boldsymbol\theta(\Omega))$, 
we obtain equivalent criteria for a contractive module map to admit a Schur-class lift: one in terms of the contractivity of an associated functional on a subspace of $L^1(\partial\boldsymbol\theta(\Omega))$, and another in terms of a geometric distance formula in the same $L^1$-space.
Specializing to quotient domains of the polydisc factored by imprimitive finite complex reflection groups, we obtain a commutant lifting criterion formulated in terms of inner functions.
Finally, we apply these operator-theoretic results to finite-point Nevanlinna--Pick type interpolation problems on $\boldsymbol \theta(\Omega)$.  
Since the symmetrized bidisc and the tetrablock arise as quotient domains of suitable bounded symmetric domains, these criteria apply in particular to those domains.




%Finally, we apply these operator-theoretic results to derive criteria for the Nevanlinna–Pick type interpolation problem for finite subsets of $\boldsymbol \theta(\Omega)$.

%obtaining Schur-class interpolation criteria for quotient domains.

%we establish multiple equivalent characterizations for a contractive module map to admit a Schur-class lift. These include a functional contractivity condition and a geometric distance formula in $L^1(\partial\boldsymbol\theta(\Omega)).$

%Finally, we apply these operator-theoretic results to solve the Nevanlinna–Pick type interpolation problem for finite subsets of $\boldsymbol \theta(\Omega)$, yielding criteria for Schur-class interpolation on domains such as the symmetrized bidisc and the tetrablock.


%In this paper, we study the criteria of commutant lifting and interpolation by Schur functions on quotient domain $\boldsymbol \theta(\Omega).$
%For quotient domains of the polydisc associated with the imprimitive pseudo-reflection group, we obtain a characterization of commutant lifting in terms of inner-function. These results are then applied to finite-point interpolation problems, yielding criteria for Schur-class interpolation on quotient domains, including the symmetrized bidisc and the tetrablock.

%As applications, we obtain interpolation criteria for prescribed finite data on quotient domains.
%In particular, the results yield interpolation criteria on the symmetrized bidisc and the tetrablock.


%we characterize contractive module maps on quotient modules that admit Schur-class lifts. The criteria are formulated both through contractivity of an associated functional on a subspace of $L^1(\partial\boldsymbol\theta(\Omega))$ and through an equivalent $L^1$-distance condition.

%The classical commutant lifting theorem of D. Sarason provides an operator-theoretic approach to the Nevanlinna–Pick interpolation problem. Such a correspondence, however, does not generally extend to the multivariable setting. In this paper, we investigate commutant lifting and interpolation problems on the symmetrized polydisc $\mathbb G_d$ in the spirit of Sarason’s approach. We obtain several characterizations of commutant lifting on quotient modules of the Hardy space $H^2(\mathbb G_d),$ including formulations in terms of distance estimates and inner functions on $\mathbb G_d.$ Using these lifting results, we derive criteria for interpolation by Schur-class functions on $\mathbb G_d.$ In particular, the interpolation results obtained here are new even for the symmetrized bidisc $\mathbb G_2.$
\end{abstract}


%D. Sarason solves the Nevanlinna-Pick interpolation problem by his commutant lifting result. But these no longer true in general in multivariable setting. This paper studies commutant lifting and interpolation problems in the setting of the symmetrized polydisc $\mathbb G_d$ using Sarason's idea. We present some characterization of commutant lifting on quotient module of the Hardy space $H^2(\mathbb G_d)$ on symmetrized polydisc. One of the characterization of commutant lifting provided in terms of inner functions on $\mathbb G_d.$ Leveraging this foundation, we present a comprehensive treatment of the interpolation problem on $\mathbb G_d.$ These interpolation criteria are in fact new for symmetrized bidisc $\mathbb G_2.$


\maketitle

\section{Introduction}
Interpolation by bounded analytic functions on the unit disc $\mathbb D \subseteq \mathbb C$ has long been a central topic in complex analysis and operator theory. Given distinct points $z_1,\ldots, z_p\in \mathbb D$ and scalars $w_1,\ldots, w_p \in \mathbb D,$ the classical Nevanlinna-Pick interpolation problem asks for necessary and sufficient condition for the existence of a holomorphic function $\varphi: \mathbb D \to \mathbb C$ satisfying $$\sup_{z\in \mathbb D} |\varphi(z)|\leq 1, \;\;\text{and}\;\;\varphi(z_i)=w_i\;\; \text{for all}\;\; i=1,\ldots, p.$$
The celebrated theorem of Nevanlinna and independently Pick asserts that such a function exists if and only if the associated $p\times p$ matrix, called the {\it  Pick matrix}, given by 
\beq \label{pick matrix} \left(\frac{1-w_i \bar{w}_j}{1-z_i \bar{z}_j}
\right)_{i,j=1}^p,\eeq is positive semi-definite.  

After nearly four decades, Donald Sarason introduced an operator-theoretic approach to the interpolation problem in \cite{DS1967}. His result, now known as the commutant lifting theorem, may be formulated as follows. Let $H^2(\mathbb D)$ be the Hardy space on $\mathbb D$, and let $T_{z}$ denote the multiplication operator on $H^2(\mathbb D)$ by the coordinate function $z$. A closed subspace $\mathcal{Q}$ of $H^2(\mathbb D)$ is called a quotient module if $T_{z}^* \mathcal{Q}\subseteq \mathcal{Q}.$ Let $P_{\mathcal{Q}}$ be the orthogonal projection of $H^2(\mathbb D)$ onto $\mathcal{Q}.$ 
Sarason’s theorem states that if $X$ is a bounded linear operator on $H^2(\mathbb D)$ satisfying $X P_{\mathcal{Q}} T_{z}|_{\mathcal{Q}}= P_{\mathcal{Q}} T_z |_{\mathcal{Q}}X$, then there exists a bounded holomorphic function $\varphi$ on $\mathbb D$ such that $$X=P_{\mathcal{Q}} T_{\varphi}|_{\mathcal{Q}} \;\; \text{and}\;\; \|X\|=\|\varphi\|_{\infty},$$ where $T_{\varphi}$ denotes the multiplication operator by $\varphi$ on $H^2(\mathbb D).$ As an application, Sarason recovered the Nevanlinna–Pick interpolation theorem.

The multivariable counterparts of both the interpolation and commutant lifting have since attracted considerable attention. In contrast to the disc, interpolation on higher dimensional domains is generally not governed by a single Pick-type matrix. The interpolation problem on the bidisc was solved by Agler and McCarthy in \cite{AM1999}. The same problem on the symmetrized bidisc $\mathbb G_2$ has been resolved recently in \cite{AY2017, BS2018}. 
There have also been several efforts to address interpolation and commutant lifting problems in general multivariable settings; see, for instance, \cite{BT1998, BLTT1999, BB2004}. However, many of these results apply only to restricted classes of functions or operators. More recently, a comprehensive characterization of both interpolation and commutant lifting for the polydisc $\mathbb D^d$
and the Euclidean unit ball $\mathbb B_d$ have been obtained in \cite{DeS2026, BDS2025}. 

The purpose of this paper is to develop commutant lifting and interpolation results on a broad class of domains that arise as the image of a bounded symmetric domain under a proper holomorphic map factored by a finite complex reflection group. 

Recall that a bounded domain $\Omega\subseteq \mathbb C^d$ is {\it symmetric} if, for every pair of points $\z, \z'\in \Omega,$ there exists an involution in the group of biholomorphic automorphisms $\mathrm{Aut}(\Omega)$ interchanging them. 
In their standard realization, bounded symmetric domains are circled, convex, and contain the origin. The open unit disc $\mathbb D$, the unit polydisc $\mathbb D^d$, and the Euclidean ball $\mathbb B_d$ in $\mathbb C^d$ are some well studied examples of bounded symmetric domains. For a masterful exposition on bounded symmetric domains, the reader is referred to \cite{LO1977, AJ1995}.

Given bounded domains $\Omega_1,\Omega_2\subseteq \mathbb C^d$, a holomorphic map $\pi: \Omega_1\to \Omega_2$ is said to be {\it proper} if the preimage of every compact subset of $\Omega_2$ is compact in $\Omega_1.$ 
Such a map is necessarily surjective, and it follows from \cite[Section 1.2]{RD1982} that there exists a positive integer $n$ (referred to as {\it multiplicity} of $\pi$) such that
$$\text{cardinality of } \pi^{-1}(\boldsymbol w) \begin{cases}
    =n, \qquad \boldsymbol w\in \Omega_2\setminus\pi(N(J_\pi))\; \text{and}\\
     < n, \qquad \boldsymbol w\in \pi(N(J_\pi)),
\end{cases}$$
where $J_\pi$ is the complex Jacobian determinant of $\pi$, and  $N(J_\pi)$ denotes its zero set.
We say $\pi$ is factored by {\it automorphisms} if there exists a finite subgroup $G\subseteq \mathrm{Aut}(\Omega_1)$ such that $\pi^{-1}(\pi(\z))=\{\sigma(\z): \sigma\in G\}$ for all $\z\in \Omega_1.$ 
It is known that such a group $G$ is either a group generated by complex reflections or conjugate to a complex reflection group, where a {\it complex reflection} is a finite order linear map $\sigma$ on $\mathbb C^d$ such that the rank of $I_d-\sigma$ is $1$. A group $G$ generated by finite number of complex reflections is referred as {\it finite complex reflection group}. 



%The proper holomorphic map $\pi$ is said to be factored by automorphisms if there exists a finite subgroup $G\subseteq \mathrm{Aut}(\Omega_1)$ such that $$\pi^{-1}\pi(\z)=\cup_{\sigma\in G}\{\sigma(\z)\}\qquad \text{for every}\;\z\in \Omega_1.$$ 
% It is known that such a group $G$ is either a complex reflection group or conjugate to a complex reflection group. A {\it complex reflection} on $\mathbb C^d$ is a linear map $\sigma:\mathbb C^d\to \mathbb C^d$ such that $\sigma$ has finite order in $GL(n,\mathbb C)$ and rank of $I_n-\sigma$ is $1.$


A bounded domain $\Omega$ in $\mathbb C^d$ is said to be $G$-invariant if $\sigma\cdot\z:=\sigma^{-1}(\z)\in \Omega$ for all $\sigma\in G.$ 
This induces an action on function $f:\Omega \to \mathbb C$ defined by  $(\sigma\cdot f)(\z):=f(\sigma^{-1}\cdot \z)$, and we say $f$ is {\it $G$-invariant} if $\sigma \cdot f=f$ for all $\sigma\in G.$ A fundamental characterization of finite complex reflection groups is that the ring of $G$-invariant polynomials is generated by a set of $d$ algebraically independent homogeneous polynomials (known as the {\it basic polynomials}) $\{\theta_1,\ldots, \theta_d\}$ associated to $G$ \cite[p. 282]{ST1954}. The mapping $\boldsymbol \theta=(\theta_1,\ldots, \theta_d): \Omega \to \boldsymbol \theta(\Omega)$ is proper holomorphic and factored by $G$ \cite{RD1982, TM2013}. Any other proper holomorphic map factored by $G$ factors through $\boldsymbol \theta$ via a biholomorphism.
Consequently, we consider $\boldsymbol \theta(\Omega)$ as the image of $\Omega$ under a proper holomorphic map factored by $G.$ 
Although the map $\boldsymbol \theta$ is not unique, the degrees of $\theta_j$'s are unique for $G$ up to order. Throughout this paper, $\Omega$ denotes a bounded symmetric domain, $G$ a finite complex reflection group acting on $\mathbb C^d$ under which $\Omega$ is invariant, and $\boldsymbol \theta(\Omega)$ the associated {\it quotient domain}, with typical elements denoted by $\boldsymbol \zeta=(\zeta_1,\ldots, \zeta_d)=\boldsymbol \theta(\z).$


Quotient domains have become an active subject of research in function theory and operator theory; see, for instance, \cite{BS1982, TM2013, MRZ2013, TirDS2021,  DS2021, BDGR2022,  GN2023, BBS2024}, as well as \cite{GR2024} and the references therein.
A suitable framework for Hardy spaces on quotient domains, together with the associated theory of Toeplitz operators, was developed in \cite{GR2024} (see also \cite{MRZ2013}). 
For our convenience, we recall in Section \ref{S2} those aspects of the Hardy space construction that are relevant to the present work.


Let $H^2(\boldsymbol \theta(\Omega))$ denote the Hardy space on the quotient domain $\boldsymbol \theta(\Omega),$ and let $\boldsymbol T_{\boldsymbol \zeta}=(T_{\zeta_1}, \ldots, T_{\zeta_d})$ be the commuting tuple of multiplication operators on $H^2(\boldsymbol \theta(\Omega))$ induced by the coordinate functions. 
%Then $\boldsymbol T_{\boldsymbol \zeta}$ is a commuting $d$-tuple of bounded linear operators on $H^2(\boldsymbol \theta(\Omega)).$ 
We denote by $H^\infty(\boldsymbol \theta(\Omega))$ the Banach algebra of bounded holomorphic functions on $\boldsymbol \theta(\Omega),$ endowed with the supremum norm $\|\cdot\|_{\infty, \boldsymbol \theta(\Omega)}.$ 
The elements of the closed unit ball $\mathcal{S}(\boldsymbol{\theta}(\Omega))$ of $H^\infty(\boldsymbol{\theta}(\Omega))$ are referred to as {\it Schur functions}.
%We write $\mathcal{S}(\boldsymbol{\theta}(\Omega))$ for the closed unit ball of $H^\infty(\boldsymbol{\theta}(\Omega))$ and its elements are called {\it Schur functions}. 
For any $\varphi \in H^\infty(\boldsymbol \theta(\Omega))$, the analytic Toeplitz operator $T_{\varphi}$ on $H^2(\boldsymbol \theta(\Omega))$ is the multiplication operator $$T_{\varphi}:H^2(\boldsymbol \theta(\Omega)) \to H^2(\boldsymbol \theta(\Omega)), \;\; T_\varphi f=\varphi f.$$ Moreover, $\|T_{\varphi}\|=\|\varphi\|_{\infty, \boldsymbol \theta(\Omega)}:=\sup\{|\varphi(\boldsymbol \zeta)|: \boldsymbol \zeta\in \boldsymbol \theta(\Omega)\}.$
Let $\mathcal{Q}$ be a quotient module of $H^2(\boldsymbol \theta(\Omega)),$ that is, a closed subspace invariant under $T_{\zeta_j}^*$ for each $ j=1,\ldots,d.$ Let $P_{\mathcal{Q}}$ denote the orthogonal projection of $H^2(\boldsymbol \theta(\Omega))$ onto $\mathcal{Q}.$ For any $\varphi\in H^\infty(\boldsymbol \theta(\Omega))$, the compression $S_{\varphi}:=P_{\mathcal{Q}} T_{\varphi}|_{\mathcal{Q}}$ defines a bounded operator on $\mathcal{Q}.$ 
In particular, $S_{\zeta_j}=P_{\mathcal{Q}} T_{\zeta_j}|_{\mathcal{Q}}\in \mathcal{B}(\mathcal{Q})$ for all $j=1,\ldots,d.$ It readily follows that $S_{\varphi }$ commutes with $S_{\zeta_j}$ for all $j,$ and that $\|S_{\varphi}\| \leq 1$ whenever $\varphi \in \mathcal{S}(\boldsymbol \theta(\Omega)).$
Motivated by this observation, we define a {\it module map} as any operator $X\in \mathcal{B}(\mathcal{Q})$ satisfying $XS_{\zeta_j}=S_{\zeta_j}X$ for all $j=1,\ldots,d.$ Thus every Schur function $\varphi\in \mathcal{S}(\boldsymbol \theta(\Omega))$ induces a contractive module map $S_{\varphi}$ on $\mathcal{Q}$.
Let $\mathcal{B}_1(\mathcal{Q})=\{X \in \mathcal{B}(\mathcal{Q}): \|X\|\leq 1\}.$ We say that a contractive module map $X\in \mathcal{B}_1(\mathcal{Q})$ is {\it liftable} or {\it admits a lift} if there exists a function $\varphi \in \mathcal{S}(\boldsymbol \theta(\Omega))$ such that $X=S_\varphi.$


The main objective of this paper is twofold. First, we establish several equivalent characterizations for the existence of Schur class lifts of contractive module maps on quotient modules of $H^2(\boldsymbol \theta(\Omega))$. These results are developed in Section \ref{S3}. In the special case of the quotient domain $\boldsymbol \theta(\mathbb D^d)$, arising from the polydisc $\mathbb D^d,$ we obtain a commutant lifting criterion expressed in terms of inner functions on $\boldsymbol \theta(\mathbb D^d)$ (see Theorem \ref{3rd characterization}). 

Secondly, in Section \ref{S4}, we apply these commutant lifting characterizations to interpolation problems. More precisely, for prescribed distinct points in $\boldsymbol \theta(\Omega)$ and prescribed target values in the unit disc, we derive criteria for the existence of a Schur class interpolant in $\mathcal{S}(\boldsymbol \theta(\Omega)).$ 
The symmetrized bidisc $\mathbb G_2$ and the tetrablock $\mathbb E$, which arise in connection with $\mu$-synthesis, are notable examples of quotient domains that have attracted considerable interest in geometric function theory and operator theory. The general results obtained in this paper apply in particular to these domains and therefore yield interpolation criteria on $\mathbb G_2$ and $\mathbb E.$





%Agle The interpolation problem for the bidisc $\mathbb D^2$ was solved by Agler and McCarthy in \cite{AM1999}.
% It is worth emphasizing that a characterization based on the positivity of a single Pick-type matrix cannot be expected in multivariable interpolation problems. Indeed, the interpolation criteria obtained in \cite{AM1999} are formulated in terms of the positive semi-definiteness of a suitable family of matrices of the form \eqref{pick matrix}.
% Continuing in this direction, the interpolation problem on the symmetrized bidisc $\mathbb G_2$ has been resolved in recent works \cite{AY2017, BS2018}. 

% that interchanges $\z$ and $\z'$, where $\mathrm{Aut}(\Omega)$ is the group of biholomorphic automorphism on $\Omega$. 
% It is well known that bounded symmetric domains are circular, convex, and contain the origin. The open unit disc $\mathbb D$, the unit polydisc $\mathbb D^d$, and the Euclidean ball $\mathbb B_d$ in $\mathbb C^d$ are some well studied examples of bounded symmetric domains. 
% Let $\mathbb K$ be the isotropy subgroup of $\mathrm{Aut}(\Omega)$ at origin, that is, $\mathbb K=\{\phi \in \mathrm{Aut}(\Omega): \phi(0)=0\}.$ For a masterful exposition on bounded symmetric domains, the reader is referred to \cite{AJ1995}. 

% Now, we discuss the notion of proper holomorphic maps. For two bounded domains $\Omega_1,\Omega_2\subseteq\mathbb C^d$, a holomorphic map $\pi: \Omega_1\to  \Omega_2$ is said to be {\it proper} if $\pi^{-1}(K)$ is a compact subset of $\Omega_1$ for every compact subset $K\subseteq \Omega_2$.
% A proper holomorphic map $\pi:\Omega_1\to \Omega_2$ is surjective and there exists a positive integer $m$, called multiplicity of $\pi$, such that $\pi:\Omega_1\setminus\pi^{-1}(\pi(N(J_\pi)))\to \Omega_2\setminus \pi(N(J_\pi))$ is a holomorphic covering with 
% $$\text{cardinality of } \pi^{-1}(\boldsymbol w) \begin{cases}
%     =m, \qquad w\in \Omega_2\setminus\pi(N(J_\pi))\; \text{and}\\
%      < m, \qquad \boldsymbol w\in \pi(N(J_\pi)),
% \end{cases}$$
% where $N(J_\pi)$ is the zero set of $J_\pi$, the determinant of the complex jacobian of $\pi$ (see \cite[chapter 15]{RD1980}).


% In this paper, our point of interest is the image $\boldsymbol \theta(\Omega)$ of a bounded symmetric domain $\Omega$ under a proper holomorphic mapping $\boldsymbol \theta$ that is factored by the finite complex reflection group $G$.
% Throughout this paper, $G$ denotes a finite complex reflection group acting on $\mathbb C^d$, $\Omega$ is a bounded symmetric domain invariant under $G$, and $\boldsymbol \theta$ denotes the quotient map associated with the action of $G$. 
% Further, the image $\boldsymbol \theta(\Omega)$ of $\Omega$ under the quotient map $\boldsymbol \theta$ is referred as {\it quotient domain} and a typical element in $\boldsymbol \theta(\Omega)$ will be denoted by $\boldsymbol \zeta=(\zeta_1,\ldots, \zeta_d)$, then $\boldsymbol \zeta=\boldsymbol \theta(\z)$ for some $\z \in \Omega.$


%We recall some of the relevant aspects of this Hardy space construction in Section \ref{S2} for our purposes.
% In recent years, the quotient domains have been studied both in the context of function theory and operator theory, see for instance \cite{BS1982,TirDS2021, BDGR2022, BR2014, DS2021, EZ2005, MRZ2013, TM2013, EKZ2013,BPR2012,BBS2024,GN2023}, and in particular \cite{ GR2024} and references therein.
%A suitable notion of Hardy space on quotient domains has been introduced in \cite{GR2024} that generalizes the notion of Hardy space studied in \cite{MRZ2013} for the domain symmetrized polydisc $\mathbb G_d.$ Further, a study of Brown Halmos type theorem has been carried out in \cite{GR2024}. Thus for a detailed description of Hardy space on the quotient domain, we refer \cite[Section 3]{GR2024}. We also briefly describe this in Section \ref{S2}.


%As an application of the general theory developed here, Section \ref{S4} specializes the interpolation results obtained here to these two domains and yields corresponding interpolation criteria on $\mathbb G_2$ and $\mathbb E$.


%recall the notion of Hardy space $H^2(\boldsymbol \theta(\Omega))$ on quotient domain $\boldsymbol \theta(\Omega)$ from \cite{GR2024} and describe the associated boundary value correspondence of elements in $H^2(\boldsymbol \theta(\Omega)).$ 
%In Section \ref{S3}, we obtain several equivalent characterization for the existence of Schur-class lift of contractive module maps on the quotient module $\mathcal{Q}$ of $H^2(\boldsymbol{\theta}(\Omega)).$ 
%For the domain $\boldsymbol \theta(\mathbb D^d)$, arises as the quotient of polydisc $\mathbb D^d$, a criterion of commutant lifting has been provided in terms of inner functions on $\boldsymbol \theta(\mathbb D^d)$ (see Theorem \ref{3rd characterization}).
% Secondly, as application of each characterization of commutant lifting developed in Section \ref{S3}, we derive a criterion of interpolation of some prescribed set of distinct points in $\boldsymbol \theta(\Omega)$ by a Schur class interpolant in $\mathcal{S}(\boldsymbol \theta(\Omega)).$ 

% Originating from $\mu$-synthesis, the domain symmetrized bidisc $\mathbb G_2$ and tetrablock $\mathbb E$ turn out to be a topic of intense research interest both in function theory and operator theory. It is known that both $\mathbb G_2$ and $\mathbb E$ are quotient domain of some suitable bounded symmetric domain. Thus the results in Section \ref{S4} gives criteria of interpolation on $\mathbb E$ and $\mathbb G_2$.  

% Let $H^2(\Omega)$ be the Hardy space of the bounded symmetric domain $\Omega.$ Since the inner product in $H^2(\Omega)$ is $\mathbb K$-invariant and the finite complex reflection group  $G\subseteq \mathbb K$, the map 
% $$G \ni \sigma \mapsto R_{\sigma}\in \mathcal{U}(H^2(\Omega))$$




% To describe the domain $\mathbb G_d$, we begin with the elementary symmetric functions. For each $1 \leq i\leq d$, let $\xi_i: \mathbb C^d\to \mathbb C$ denote the $i$th elementary symmetric function, given by $$ \xi_i(\z):=\sum_{1 \leq j_1< j_2 < \cdots < j_i\leq d} z_{j_1}\cdots z_{j_i},\qquad \z=(z_1,\ldots,z_d)\in \mathbb C^d.$$
%  The symmetrization map $\\xi:\mathbb C^d \to \mathbb C^d$, defined by $$\\xi(\z)=\left(\xi_1(\z),\ldots, \xi_d(\z)\right),\qquad \z\in \mathbb C^d,$$ is a proper holomorphic map with multiplicity $d!$ \cite{RD1982}. The symmetrized polydisc $\mathbb G_d$ is then defined by $\mathbb G_d:=\\xi(\mathbb D^d)$; it is known to be a non-convex yet polynomially convex domain, see \cite{EZ2005}.
% Throughout this paper, a typical point in $ \mathbb G_d$ will be denoted by $\\xi=(\xi_1,\ldots,\xi_d).$
% It is worth noting that the symmetrized polydisc is a higher dimensional analogue of the symmetrized bidisc $\mathbb G_2:=\{(z_1+z_2,z_1z_2): z_1, z_2 \in \mathbb D\},$ which arises naturally in $\mu$-synthesis and has been extensively studied in complex geometry, function theory, and operator theory \cite{AY2004, BS2018, ALY2019, TirDS2021}.
% In recent years, the higher-dimensional analog $\mathbb G_d$ has also attracted significant attention, particularly in relation to its complex geometry \cite{EZ2005, NN2006, NPZ2007, NPTZ2008, GS2016}, function theory \cite{TTZ2011, MRZ2013, NTT2016}, and operator theory \cite{BR2014, DS2021}.
% For example, using the transformation rule for the Bergman kernel under proper holomorphic maps \cite[Theorem 1]{BS1982}, an explicit formula for the Bergman kernel on $\mathbb G_d$ was obtained in \cite{EZ2005}. Subsequently, weighted Bergman spaces and Hardy space on $\mathbb G_d,$ along with their reproducing kernel, were studied in \cite{MRZ2013}. 






% We now briefly outline the structure of the paper. 

% In Section \ref{S2}, we collect the necessary preliminaries on the Hardy space $H^2(\mathbb G_d)$ and related structures that will be used throughout the paper. 
% In particular, we recall the definition of $H^2(\mathbb G_d)$ and the associated space $L^2(\partial\mathbb G_d)$ of square-integrable functions on the Shilov boundary $\partial\mathbb G_d$ with respect to the measure $\mu_{\\xi}.$ 
% In analogy with the polydisc setting, we establish an isometric embedding of $H^2(\mathbb G_d)$ into $L^2(\partial\mathbb G_d),$ which associates to each function in $H^2(\mathbb G_d)$ an appropriate boundary function without altering its norm. 
% Furthermore, we describe orthonormal bases for both $H^2(\mathbb G_d)$ and $L^2(\partial\mathbb G_d),$ and discuss their relationship with the classical spaces $H^2(\mathbb D^d)$ and $L^2(\mathbb T^d).$

% Section \ref{S3} is devoted to the study of commutant lifting for contractive module maps on quotient modules of $H^2(\mathbb G_d).$ For a given quotient module $\mathcal{Q},$ we obtain several equivalent characterizations for the existence of lifts of contractive module maps $X\in \mathcal{B}_1(\mathcal{Q}).$
% The methods and results are inspired by recent developments in the polydisc setting \cite{DeS2026}. In particular, one of our characterizations is formulated in terms of inner functions on $\mathbb G_d.$ 
% As a key step, we show that every Schur function $\varphi\in \mathcal{S}(\mathbb G_d)$ arises as a weak* limit of a sequence of inner functions on $\mathbb G_d,$ by suitably adapting the corresponding result for the polydisc $\mathbb D^d.$

% In Section \ref{S4}, we address an interpolation problem on the symmetrized polydisc. Given distinct points $\\xi_1,\\xi_2,\ldots, \\xi_m \in \mathbb G_d$ and prescribed values $w_1,\ldots, w_d\in \mathbb D$, we seek necessary and sufficient conditions for the existence of a function $\varphi\in \mathbb G_d$ such that $$\varphi(\boldsymbol \xi_i)=w_i, \; i=1,2,\ldots,m.$$ 
% By employing the commutant lifting criteria developed in Section \ref{S3}, we obtain a complete characterization of this interpolation problem.




\section{Preliminaries}\label{S2}
In this section, we recall the Hardy space on the quotient domain $\boldsymbol \theta(\Omega)$ and describe its boundary realization. Recall that $\boldsymbol \theta=(\theta_1,\ldots, \theta_d):\Omega \to \boldsymbol \theta(\Omega)$ is the proper holomorphic map factored by the finite complex reflection group $G.$ We write $\partial\Omega$ and  $\partial\boldsymbol \theta(\Omega)$ for the Shilov boundaries of $\Omega$ and $\boldsymbol \theta(\Omega)$, respectively.
%We denote the Shilov boundaries of $\Omega$ and $\boldsymbol \theta(\Omega)$ by $\partial\Omega$ and  $\partial\boldsymbol \theta(\Omega)$, respectively. 
Since the quotient map $\boldsymbol\theta$ extends as a proper holomorphic map of the same multiplicity to suitable neighborhoods of the corresponding closures, it follows from \cite[p. 100, Corollary 3.2]{KZ2013} that $\partial\boldsymbol \theta(\Omega)=\boldsymbol \theta(\partial\Omega).$ 

% The quotient map $\boldsymbol \theta:\Omega \to \boldsymbol \theta(\Omega)$ can be extended to a proper holomorphic map of the same multiplicity from a neighborhood of $\Omega$ onto a neighborhood of $\boldsymbol \theta(\Omega).$
% Then \cite[p. 100, Corollary 3.2]{KZ2013} states that $\boldsymbol \theta ^{-1}(\partial\boldsymbol \theta(\Omega))=\partial\Omega$. Thus the Shilov boundary $\partial\boldsymbol \theta(\Omega)$ of $\boldsymbol \theta(\Omega)$ satisfies $\partial\boldsymbol \theta(\Omega)=\boldsymbol \theta(\partial\Omega).$ 


Let $\mathbb K=\{\phi\in \mathrm{Aut}(\Omega): \phi(0)=0\}$ be the isotropy subgroup at the origin, and let $\mu$ denote the normalized $\mathbb K$-invariant Lebesgue measure on $\partial\Omega.$ Then the Hardy space $H^2(\Omega)$ consists of all holomorphic functions $f:\Omega\to \mathbb C$ such that  
$$\|f\|_{H^2(\Omega)}^2=\sup_{0<r<1} \int_{\partial\Omega} |f(r\z)|^2d\mu(z) <\infty.$$
Let $L^2(\partial\Omega)$ denote the space of square-integrable functions with respect to $\mu.$ By \cite[Theorem 6]{HM1969}, every $f\in H^2(\Omega)$ has a boundary value function $f^* \in L^2(\partial\Omega)$ satisfying $$\lim_{r \to 1^-} \int_{\partial\Omega}|f_r-f^*|^2 d\mu =0,\quad \text{and} \quad \|f\|_{H^2(\Omega)}=\|f^*\|_{L^2(\partial\Omega)},$$ where the radial slice  $f_r:\partial\Omega \to \mathbb C$ is given by $f_r(\z)=f(r\z)$ for $0<r<1$ (see also \cite[p. 126]{UH1996}). Thus, the boundary value map $i: H^2(\Omega) \to L^2(\partial\Omega)$, defined by $i(f)=f^*$, embeds $H^2(\Omega)$ isometrically into $L^2(\partial\Omega)$. We denote its range by $H^2(\partial\Omega)$. 
% This leads to an isometric embedding $i: H^2(\Omega) \to L^2(\partial\Omega)$ of $H^2(\Omega)$ onto a closed subspace $H^2(\partial\Omega)$ of $L^2(\partial\Omega).$

The space $H^2(\Omega)$ is a reproducing kernel Hilbert space. Its reproducing kernel, denoted by $\mathbb S_\Omega$, is the Szeg\"o kernel of $\Omega.$
The corresponding Poisson kernel $P:\Omega \times \partial\Omega  \to \mathbb C$ is given by $$P( \z, \boldsymbol \xi)=\frac{|\mathbb S_\Omega(\z, \boldsymbol \xi)|^2}{\mathbb S_\Omega(\z,\z)}, \qquad \z \in \Omega,\; \boldsymbol \xi \in \partial\Omega.$$ By \cite[Theorem 1]{HM1973} (see also \cite{KA1965}), any $f\in H^2(\Omega)$ can be recovered from its boundary value $f^* \in L^2(\partial\Omega)$ via the Poisson integral formula:
$$f(\z)=\int_{\partial\Omega}  P(\z, \boldsymbol \xi) f^*(\boldsymbol \xi)d\mu(\boldsymbol \xi), \qquad \z \in \Omega.$$
Moreover, $f_r$ converges to $f^*$ almost everywhere on $\partial\Omega$ (see \cite[Theorem 2]{WN1968}), that is,
$$f^*(\z)=\lim_{r \to 1^-} f(r\z)\qquad \text{for a.e}\; \z \in \partial \Omega.$$  

We now recall the Hardy space on the quotient domain. Using the change of variables formula, define a measure $\mu_{\boldsymbol\theta}$ supported on $\partial \boldsymbol \theta(\Omega) $ by
$$\int_{\partial\boldsymbol{\theta} (\Omega)} f d\mu_{\boldsymbol\theta}= \int_{\partial\Omega} (f\circ \boldsymbol\theta) |J_{\boldsymbol \theta}|^2 d\mu,$$ 
where $J_{\boldsymbol\theta}$ denotes the complex Jacobian determinant of the proper holomorphic map $\boldsymbol\theta: \Omega \to \boldsymbol \theta (\Omega).$ The Hardy space $H^2(\boldsymbol \theta(\Omega))$  is then defined as the Hilbert space of holomorphic functions $f$ on $\boldsymbol \theta(\Omega)$ for which $J_{\boldsymbol \theta} (f\circ \boldsymbol \theta)\in H^2(\Omega)$. Furthermore, the norm of a function $f\in H^2(\boldsymbol \theta(\Omega))$ is given by 
\beq \label{eq 4} \|f\|_{H^2(\boldsymbol \theta(\Omega))}=\frac{1}{c_{\boldsymbol \theta}} \left\{\sup_{0<r<1} \int_{\partial\Omega}|f\circ \boldsymbol\theta(r\z)|^2 |J_{\boldsymbol\theta}(r\z)|^2 d\mu(\z) \right\}^{1/2},\eeq where the normalization by $c_{\boldsymbol \theta}=\|J_{\boldsymbol\theta}\|_{H^2(\Omega)}$ ensures $\|1\|_{H^2(\boldsymbol \theta(\Omega))}=1.$
It follows from \cite[Proposition 3.13]{GR2024} that polynomials are dense in $H^2(\boldsymbol \theta(\Omega))$ and that it is a reproducing kernel Hilbert space with the reproducing kernel $\mathbb S_{\boldsymbol \theta(\Omega)}:\boldsymbol \theta(\Omega) \times \boldsymbol \theta(\Omega) \to \mathbb C$ given by 
\beq \label{Hardy kernel} \mathbb S_{\boldsymbol \theta(\Omega)}(\boldsymbol \theta(\z), \boldsymbol\theta(\boldsymbol w))=\frac{c_{\boldsymbol \theta}^2}{|G|}\frac{1}{J_{\boldsymbol \theta}(\z)\overline{J_{\boldsymbol \theta}(\boldsymbol w)}}  \sum_{\sigma\in G} \sgn(\sigma^{-1}) \mathbb S_{\Omega}(\sigma^{-1}\cdot \z, \boldsymbol w),\qquad \z, \boldsymbol w\in \Omega.\eeq
%(see \cite[equation 3.4]{GR2024}.) 

Let $L^2(\partial\boldsymbol \theta(\Omega))$ denote the space of square-integrable functions with respect to $\mu_{\boldsymbol\theta}.$ For $f\in L^2(\partial\boldsymbol \theta(\Omega)),$ the norm is given by 
\beq \label{eq 5} \|f\|_{L^2(\partial\boldsymbol \theta(\Omega))}= \frac{1}{c_{\boldsymbol \theta}} \|J_{\boldsymbol \theta} (f\circ \boldsymbol \theta)\|_{L^2(\partial \Omega)}= \frac{1}{c_{\boldsymbol \theta}}\left( \int_{\partial\Omega} |f(\boldsymbol \theta(\z))|^2 |J_{\boldsymbol \theta}(\z)|^2 d\mu(\z)\right)^{1/2}.\eeq

Recall that the proper map $\boldsymbol \theta=(\theta_1,\ldots, \theta_d) :\Omega \to \boldsymbol \theta(\Omega)$ consists of homogeneous polynomials; let $k_j$ denote the degree of $\theta_j$ for each $1\leq j\leq d.$  
If $f\in H^2(\boldsymbol \theta(\Omega))$, then $J_{\boldsymbol \theta}(f \circ \boldsymbol \theta) \in H^2(\Omega)$, and its boundary value $(J_{\boldsymbol \theta}(f \circ \boldsymbol \theta) )^*$ belongs to $ H^2(\partial\Omega).$ Since $J_{\boldsymbol \theta}$ is a homogeneous polynomial, the following limit exists for almost every $\z\in \partial\Omega$: \beq \label{eq 16}(J_{\boldsymbol \theta}(f \circ \boldsymbol \theta) )^*(z)=\lim_{r \to 1^-} J_{\boldsymbol \theta}(r\z)(f \circ \boldsymbol \theta)(r\z)=J_{\boldsymbol \theta}(\z) \lim_{r \to 1^-}f(r^{k_1} \theta_1(\z), \ldots, r^{k_d}\theta_d(\z)).\eeq
Let $V\subseteq \partial\Omega$ be the set of points where the limit in \eqref{eq 16} does not exist. Then the set $V$ is $G$-invariant and $\mu(V)=0$. Since the zero set $N(J_{\boldsymbol \theta})$ of $J_{\boldsymbol \theta}$ is $G$-invariant with $\mu(N(J_{\boldsymbol \theta}))=0$, we define a function $\tilde{f}^*$ on $\partial\boldsymbol \theta(\Omega)$ by
$$\tilde{f}^* (\boldsymbol \zeta)=\begin{cases}
    \frac{(J_{\boldsymbol \theta}(f \circ \boldsymbol \theta) )^*(z)}{J_{\boldsymbol \theta}(\z)}, \; \text{ if } \boldsymbol \zeta=\boldsymbol \theta(\z)\in \partial\boldsymbol \theta(\Omega) \text{ with } \z \notin V\cup N(J_{\boldsymbol \theta})\\
    0, \qquad\qquad\;\; \text{otherwise}
\end{cases}.$$
%where $N(J_{\boldsymbol \theta})$ is the zero set on $\partial\Omega$ of $J_{\boldsymbol \theta}.$ Note that $\mu(N(J_{\boldsymbol \theta}))=0.$
This gives a well defined element of $L^2(\partial\boldsymbol \theta(\Omega)).$ 
%definition is independent of the choice of $\z$ for which $\boldsymbol \zeta=\boldsymbol \theta(\z).$
Given a function $f:\boldsymbol \theta(\Omega) \to \mathbb C$ and $0<r<1$, we define the slice $\tilde{f}_r:\partial\boldsymbol \theta(\Omega)\to \mathbb C$ by $$\tilde{f}_r(\boldsymbol \zeta)=f(r^{k_1}\zeta_1, r^{k_2} \zeta_2,\ldots, r^{k_d} \zeta_d),\qquad \boldsymbol \zeta=(\zeta_1,\ldots,\zeta_d)\in \partial\boldsymbol \theta(\Omega).$$ 
If $f\in H^2(\boldsymbol \theta(\Omega)),$ a routine verification shows that  $\tilde{f}_r \in L^2(\partial\boldsymbol \theta(\Omega))$ for all $0<r<1$, and that $\tilde{f}^*$ is measurable and square integrable with respect to $\mu_{\boldsymbol \theta}$, satisfying $$\lim_{r \to 1^-} \int_{\partial\boldsymbol \theta(\Omega)} |\tilde{f}_r-\tilde{f}^*|^2d\mu_{\boldsymbol \theta}=0.$$ Moreover, for all $\boldsymbol \zeta=\boldsymbol\theta(\z)$ with $\z\notin V\cup N(J_{\boldsymbol \theta})$, we have  $\tilde{ f}^*(\boldsymbol \zeta)=\lim_{r\to 1^-} \tilde{f}_r(\boldsymbol \zeta)$, and $J_{\boldsymbol \theta} (\tilde{f}^* \circ \boldsymbol \theta) \in H^2(\partial\Omega).$




Define $H^2(\partial\boldsymbol \theta(\Omega)):=\{f\in  L^2(\partial\boldsymbol \theta(\Omega)): J_{\boldsymbol \theta}(f \circ \boldsymbol \theta) \in H^2(\partial\Omega)\}.$ 
%Then $H^2(\partial\boldsymbol \theta(\Omega))$ is a closed subspace of $L^2(\partial\boldsymbol \theta(\Omega))$. 
%denote the closed subspace of $L^2(\partial\boldsymbol \theta(\Omega))$ consisting of those functions $f\in L^2(\partial\boldsymbol \theta(\Omega))$ such that $J_{\boldsymbol \theta}(f \circ \boldsymbol \theta) \in H^2(\partial\Omega).$ 
The following result establishes an isometric embedding of $H^2(\boldsymbol \theta(\Omega))$ onto the closed subspace $H^2(\partial\boldsymbol \theta(\Omega))$ of $L^2(\partial\boldsymbol \theta(\Omega)).$

\begin{proposition}\label{required result 1}
    Let $i': H^2(\boldsymbol \theta(\Omega)) \to L^2(\partial\boldsymbol {\theta}(\Omega))$ be the map given by $i'(f)=\tilde{f}^*.$ Then $i'$ is an isometry, and its range is precisely $H^2(\partial\boldsymbol \theta(\Omega)).$
\end{proposition}

\begin{proof} Let $\sgn(\cdot)$ denote the character of the sign representation of finite complex reflection group $G.$ A function $f$ in $d$ complex variables is said to be {\it $G$-alternating} if $\sigma\cdot f=\sgn(\sigma) f$ for all $\sigma\in G.$
  To relate the spaces $H^2(\boldsymbol \theta(\Omega))$ and $L^2(\partial\boldsymbol \theta(\Omega))$ to their corresponding function spaces on $\Omega$ and $\partial\Omega$, we consider the following subspaces:
\begin{align*}
    & H_{\sgn}^2(\Omega):=\left\{f\in H^2(\Omega): f \text{ is } G\text{-alternating}\right\}, \text{ and}\\
& L^2_{\sgn}(\partial\Omega):=\left\{h\in L^2(\partial\Omega): h \text{ is } G\text{-alternating }\mu \text{-a.e.} \right\}.
\end{align*} 
Clearly, the boundary value map $i: H^2(\Omega) \to H^2(\partial\Omega)$, when restricted to $H^2_{\sgn}(\Omega)$, surjects onto $H^2_{\sgn}(\partial\Omega):= L^2_{\sgn}(\partial\Omega) \cap H^2(\partial\Omega).$ 
 To transport this structure to the quotient domain $\boldsymbol \theta(\Omega)$, we consider the operators $\Gamma_h:H^2(\boldsymbol \theta(\Omega)) \to H^2_{\sgn}(\Omega)$ and  $\Gamma_{l}: L^2(\partial\boldsymbol \theta(\Omega))\to L_{\sgn}^2(\partial\Omega)$, defined by $$\Gamma_h f=\frac{1}{c_{\boldsymbol \theta}} J_{\boldsymbol \theta} (f \circ \boldsymbol \theta), \; f\in H^2(\boldsymbol \theta(\Omega)) \; \text{ and }\;  \Gamma_{l} g=\frac{1}{c_{\boldsymbol \theta}} J_{\boldsymbol\theta} (g \circ \boldsymbol\theta),\; g\in L^2(\partial\boldsymbol \theta(\Omega)).$$  
In view of \eqref{eq 4} and \eqref{eq 5}, both $\Gamma_h$ and $\Gamma_l$ are (well defined) isometries. In fact, they are unitary operators (see \cite[Lemma 3.5, Lemma 3.14]{GR2024}). This yields the following commutative diagram:
$$\begin{tikzcd}
H^2(\boldsymbol \theta(\Omega))\arrow{r}{\Gamma_l^{*} \circ i \circ \Gamma_h } \arrow[swap]{d}{\Gamma_h} &  L^2(\partial \boldsymbol \theta(\Omega)) \arrow{d}{\Gamma_l} \\
 H^2_{\sgn}(\Omega) \arrow{r}{i} &L^2_{\sgn}(\partial\Omega)
\end{tikzcd}.$$
 Hence, $\Gamma_l^{*} \circ i \circ \Gamma_h: H^2(\boldsymbol{\theta}(\Omega))\to L^2(\partial\boldsymbol \theta(\Omega))$ is an isometric embedding of $H^2(\boldsymbol \theta(\Omega))$ onto a closed subspace of $L^2(\partial\boldsymbol \theta(\Omega))$ (see \cite[Lemma 3.15]{GR2024}). It remains to verify that $i'=\Gamma_l^{*} \circ i \circ \Gamma_h$ and that its range is $H^2(\partial\boldsymbol \theta(\Omega)).$  
 
Let $f\in H^2(\boldsymbol \theta(\Omega)).$ By the preceding discussion, $\tilde{f}^*\in L^2(\partial\boldsymbol \theta(\Omega))$ and $J_{\boldsymbol \theta}(\tilde{f}^* \circ \theta)\in H_{\sgn}^2(\partial\Omega).$ Moreover, for almost every $\z\in \partial\Omega$, $$i(\Gamma_h(f))(\z)=\lim_{r \to 1^-} J_{\boldsymbol \theta}(r\z) f(\boldsymbol{\theta}(r\z))=J_{\boldsymbol \theta}(\z)\lim_{r\to 1^-} \tilde{f}_r(\boldsymbol{\theta}(\z)) =J_{\boldsymbol \theta}(\z) \tilde{f}^*(\boldsymbol{\theta}(\z)).$$
Consequently, $\Gamma_l^* \circ i \circ \Gamma_h (f)= \Gamma_l^*(J_{\boldsymbol{\theta}} \tilde{f}^* \circ \boldsymbol{\theta})=\tilde{f}^*=i'(f).$

To establish surjectivity, let $\tilde{f}^*\in H^2(\partial\boldsymbol \theta(\Omega)).$  
Equivalently,  $$\Gamma_l(\tilde{f}^*)=\frac{1}{c_{\boldsymbol \theta}} J_{\boldsymbol \theta} \tilde{f}^* \circ \boldsymbol \theta\in H^2_{\sgn}(\partial\Omega)=i(H^2_{\sgn}(\Omega)).$$ Hence, there exists $g\in H^2_{\sgn}(\Omega)$ such that $\frac{1}{c_{\boldsymbol \theta}} J_{\boldsymbol\theta} \tilde{f}^* \circ \boldsymbol\theta=g^*=i(g)\in H^2_{\sgn}(\partial\Omega).$ Setting $f=\Gamma_h^*(g)\in H^2(\boldsymbol \theta(\Omega)),$ a straightforward verification shows that $i'(f)=\tilde{f}^*.$ Therefore, $i'(H^2(\boldsymbol \theta(\Omega)))=H^2(\partial\boldsymbol \theta(\Omega)).$
\end{proof}

Let $L^\infty(\partial\boldsymbol \theta(\Omega))$ denote the algebra of essentially bounded $\mu_{\boldsymbol \theta}$-measurable functions on $\partial\boldsymbol \theta(\Omega)$, and let 
%The space of all $G$-invariant essentially bounded measurable functions on $\partial\Omega$ is denoted by 
$L^\infty(\partial\Omega)^{G}=\{ f\in L^\infty (\partial\Omega): \sigma(f)=f\;  \mu\text{-a.e. for all } \sigma\in G\}.$ There exists an isometric $*$-isomorphism from $L^\infty(\partial\boldsymbol \theta(\Omega))$ onto $L^\infty(\partial\Omega)^{G}$ (see \cite[Subsection 3.6]{GR2024}). 
We define $$H^\infty(\partial \boldsymbol \theta(\Omega)):=H^2(\partial\boldsymbol \theta(\Omega)) \cap L^\infty(\partial\boldsymbol \theta(\Omega)).$$ Note that $H^\infty(\boldsymbol \theta(\Omega)) \subseteq H^2(\boldsymbol \theta(\Omega)).$ The following result shows that the boundary value map
$i'$ restricts to an isometric isomorphism between $H^\infty(\boldsymbol \theta(\Omega))$ and $H^\infty(\partial\boldsymbol \theta(\Omega)).$

\begin{proposition} \label{required result 2}
    For every $f\in H^\infty(\boldsymbol\theta(\Omega))$, the boundary value function $\tilde{f}^* $ belongs to $H^\infty(\partial\boldsymbol \theta(\Omega))$ and satisfies $\|\tilde{f}^*\|_{L^\infty(\partial\boldsymbol \theta(\Omega))}=\|f\|_{\infty, \boldsymbol\theta(\Omega)}.$ Furthermore, $i'(H^\infty(\boldsymbol \theta(\Omega)))=H^\infty(\partial\boldsymbol \theta(\Omega)).$ 
\end{proposition}
\begin{proof}
    Let $f\in H^\infty(\boldsymbol\theta(\Omega)).$ Then $\tilde{f}^* \in H^2(\partial\boldsymbol \theta(\Omega)).$ Since $|\tilde{f}_r(\boldsymbol \zeta)| \leq \|f\|_{\infty, \boldsymbol\theta(\Omega)}$ for all $0<r<1$ and $\boldsymbol \zeta \in \partial\boldsymbol \theta(\Omega),$ passing to the radial limits almost everywhere yields $\|\tilde{f}^*\|_{L^\infty(\partial\boldsymbol \theta(\Omega))} \leq \|f\|_{\infty, \boldsymbol \theta(\Omega)}.$ This confirms that $\tilde{f}^* \in H^\infty(\partial\boldsymbol \theta(\Omega)).$

    Conversely, suppose $h\in H^\infty (\partial \boldsymbol \theta(\Omega)).$ By Proposition \ref{required result 1}, there exists $f\in H^2(\boldsymbol \theta(\Omega))$ such that $i'(f)=h.$ In particular, $$\lim_{r \to 1^-} \tilde{f}_r(\boldsymbol \zeta)=h(\boldsymbol \zeta)\qquad \text{for almost every } \boldsymbol \zeta \in \partial\boldsymbol \theta(\Omega).$$ Consequently, for any polynomial $p$ in $d$-variables, $i'(pf)=p|_{\partial\boldsymbol \theta(\Omega)} h.$
    We show that $f$ is bounded and that its boundary function is $h.$
    Since $h \in L^\infty(\partial\boldsymbol \theta(\Omega))$, the multiplication operator $M_{h}: L^2(\partial\boldsymbol \theta(\Omega)) \to L^2(\partial\boldsymbol \theta(\Omega))$ defined by $M_{h}(g)= h g$ is bounded, with $\|M_{h}\|=\|h\|_{L^\infty(\partial\boldsymbol \theta(\Omega))}.$
    For every polynomial $p$, $$M_{h}(p|_{\partial\boldsymbol \theta(\Omega)})=p|_{\partial\boldsymbol \theta(\Omega)} h= i'(pf)\in H^2(\partial\boldsymbol\theta(\Omega)).$$ 
    As polynomials are dense in $H^2(\partial\boldsymbol \theta(\Omega))$, it follows that $M_{h}(H^2(\partial\boldsymbol \theta(\Omega))) \subseteq H^2(\partial\boldsymbol \theta(\Omega)).$
    %Now, we observe that $M_{\tilde{f}^*}(H^2(\partial\boldsymbol \theta(\Omega))) \subseteq H^2(\partial\boldsymbol \theta(\Omega)).$ 
    Consequently, $T:=i'^{-1} \circ M_{h}\circ i'$ is a bounded operator on $H^2(\boldsymbol \theta(\Omega))$, and $\|T\| \leq \|h\|_{L^\infty(\partial\boldsymbol \theta(\Omega))}.$
   % Now, we show that $f$ is a multiplier on the reproducing kernel Hilbert space $H^2(\boldsymbol \theta(\Omega)).$ 
    %To this begin, consider $T:=i'^{-1} \circ M_{\tilde{f}^*}\circ i' \in \mathcal{B}(H^2(\boldsymbol \theta\Omega)),$ then $\|T\| \leq \|\tilde{f}^*\|_{L^\infty(\partial\boldsymbol \theta(\Omega))}.$
    Moreover, for any polynomial $p$, $$fp=i'^{-1}( p|_{\partial\boldsymbol\theta(\Omega)}h)=i'^{-1}(M_{h} (i'(p)))=Tp\in H^2(\boldsymbol \theta(\Omega)).$$ 
    %Since polynomials are dense in $H^2(\boldsymbol \theta(\Omega))$, 
    Hence, $f$ is a multiplier on the reproducing kernel Hilbert space $H^2(\boldsymbol \theta(\Omega))$, and the corresponding multiplier operator $\mathcal{M}_f$ is precisely $T$, ensuring $\|\mathcal{M}_f\| \leq \|h\|_{L^\infty(\partial\boldsymbol \theta(\Omega))}.$ 
    By the reproducing property, for any $\boldsymbol \zeta\in \boldsymbol \theta(\Omega)$, we have $\mathcal{M}_f^*(\mathbb S_{\boldsymbol \theta(\Omega)}(\cdot, \boldsymbol{\zeta}))=\overline{f(\boldsymbol \zeta)} \mathbb S_{\boldsymbol \theta(\Omega)}(\cdot, \boldsymbol{\zeta}),$ where $\mathbb S_{\boldsymbol \theta(\Omega)}(\cdot, \boldsymbol{\zeta})\in H^2(\boldsymbol \theta(\Omega))$ is the reproducing kernel at $\boldsymbol \zeta$, given by $\mathbb S_{\boldsymbol \theta(\Omega)}(\cdot, \boldsymbol{\zeta})(\boldsymbol \eta)=\mathbb S_{\boldsymbol \theta(\Omega)}(\boldsymbol \eta, \boldsymbol{\zeta})$ for all $\boldsymbol \eta\in \boldsymbol \theta(\Omega).$ 
    Therefore, $$|f(\boldsymbol \zeta)| \leq \frac{\|\mathcal{M}_f^*(\mathbb S_{\boldsymbol \theta(\Omega)}(\cdot, \boldsymbol{\zeta}))\|}{\|\mathbb S_{\boldsymbol \theta(\Omega)}(\cdot, \boldsymbol{\zeta})\|}\leq \|\mathcal{M}_f\| \leq \|h\|_{L^\infty(\partial\boldsymbol \theta(\Omega))}.$$ 
    Thus, $f$ is bounded and holomorphic on $\boldsymbol \theta (\Omega)$ with $\|f\|_{\infty, \boldsymbol \theta(\Omega)} \leq \|h\|_{L^\infty(\partial\boldsymbol \theta(\Omega))}.$
    On the other hand, the fact that $\lim_{r \to 1^-} \tilde{f}_r(\boldsymbol \zeta)=h(\boldsymbol \zeta)$  almost everywhere on $\boldsymbol \zeta \in \partial\boldsymbol \theta(\Omega)$ yields that $\|h\|_{L^\infty(\partial\boldsymbol \theta(\Omega))} \leq \|f\|_{\infty, \boldsymbol \theta(\Omega)}.$
    This establishes the isometric equality and completes the proof.
    \end{proof}





% The sign representation of a finite complex reflection group $G$, $\sgn : G\to \mathbb C^*$, is defined by \cite[p. 139, Remark (1)]{RS1977} $\sgn(\sigma)=(\det(\sigma))^{-1}$ for all $\sigma\in G.$
% A function $f:\Omega \to \mathbb C$ is said to be {\it $G$-invariant} if $\sigma\cdot f=f$ for all $\sigma\in G$, and $f$ is said to be {\it $G$-alternating} if $\sigma\cdot f=\sgn(\sigma) f$ for all $\sigma\in G.$
% In this section, we collect preliminary material on the Hardy space $H^2(\boldsymbol \theta(\Omega))$ of proper image $\boldsymbol \theta(\Omega)$ and its boundary realization. We begin by recalling  that $\partial\Omega$ and  $\partial\boldsymbol \theta(\Omega)$ are the Shilov boundaries of $\Omega$ and  $\boldsymbol \theta(\Omega),$ respectively, and that the proper holomorphic map $\boldsymbol \theta$ satisfies $\boldsymbol\theta(\partial\Omega)=\partial\boldsymbol \theta(\Omega).$

%In order to discuss the Hardy space on quotient domains, we recall that $\sgn : G\to \mathbb C^*$, the {\it sign representation} of the finite complex reflection group $G$, is given by $\sgn(\sigma)=(\det(\sigma))^{-1}$ for all $\sigma\in G$ \cite[p. 139, Remark (1)]{RS1977}.

%Since $J_{\boldsymbol \theta}(\z)$ is a homogeneous polynomial of some fix degree $k$, we have  $J_{\boldsymbol \theta}(r\z)=r^k J_{\boldsymbol \theta}(\z).$ 

% Moreover, since $J_{\boldsymbol \theta} f\circ \boldsymbol \theta\in H^2(\Omega)$ and the zero set of $J_{\boldsymbol{\theta}}$ on $\partial\Omega$ has measure zero, the existence of the radial limit 
% $$\lim_{r \to 1^-} J_{\boldsymbol \theta}(r\z) f(\boldsymbol{\theta}(r\z))=\lim_{r\to 1^-} r^k J_{\boldsymbol\theta} (\z) \tilde{f}_r(\theta_1(\z),\ldots, \theta_d(\z))$$ implies that $\lim_{r \to 1^-}\tilde{f}_r(\q)$ exists for $\mu_{\boldsymbol \theta}$-a.e. $\q\in \partial\boldsymbol\theta(\Omega).$ Consequently, $\tilde{f}^* \in L^2(\partial\boldsymbol \theta(\Omega)),$ and $\tilde{f}_r \to \tilde{f}^*$ in $L^2(\partial\boldsymbol \theta(\Omega)).$

%A function $f:\Omega \to \mathbb C$ is said to be {$G$-alternating} if $\sigma\cdot f=\sgn(\sigma) f$ for all $\sigma\in G.$ Since the Jacobian $J_{\boldsymbol \theta}$ is $G$-alternating, so is  $ (J_{\boldsymbol \theta}(f \circ \boldsymbol \theta)$ for all $f\in H^2(\Omega)$.

% Recall that for any $\sigma\in G$ and $\z=(z_1,\ldots, z_d)\in \mathbb C^d$, $\sigma\cdot \z=\sigma^{-1}(\z).$ For any function $f$ on $\boldsymbol \theta(\Omega)$, the group $G$ acts on $f$ by $\sigma(f)(\z)= f( \sigma^{-1}\cdot \z).$ 


%that is, $J_{\boldsymbol\theta} (\tilde{f}^* \circ \boldsymbol\theta)\in H^2_{\sgn}(\partial\Omega).$
    

%     Conversely, suppose $\tilde{f}^* \in H^\infty(\partial\boldsymbol \theta(\Omega)).$ Then $\tilde{f}^* \circ \boldsymbol \theta \in L^\infty(\partial\Omega)$ and $g^*=J_{\boldsymbol \theta}(\tilde{f}^* \circ \boldsymbol \theta)\in H^2(\partial\Omega)$. Since $J_{\boldsymbol \theta}$ is bounded on $\partial\Omega$, $J_{\boldsymbol \theta}(\tilde{f}^* \circ \boldsymbol \theta) \in L^{\infty}(\partial\Omega)$, that is, $g^* \in H^\infty(\partial\Omega)=H^2(\partial\Omega)\cap L^\infty(\partial\Omega).$ Thus there exists a bounded holomorphic function $g\in H^\infty(\Omega)$ such that $i(g)=g^*$ with $\|g\|_{\infty, \Omega}=\|g\|_{L^\infty(\partial\Omega)}.$ Moreover, $g$ is anti-symmetric, that is, $\sigma(g)=\sgn(\sigma) g$ for all $\sigma\in G.$ 
% %     By Proposition \ref{required result 1}, there exists $f\in H^2(\boldsymbol \theta(\Omega))$ such that $i'(f)=\tilde{f}^*.$ Then for almost every $\q\in \partial\boldsymbol \theta(\Omega)$, $\lim_{r\to 1^-} \tilde{f}_r(\q)=\tilde{f}^*(\q).$ Defining $g=J_{\boldsymbol\theta} (f\circ \boldsymbol\theta)$, we have $g^*(\z)=\lim_{r\to 1^-} g_r(\z)= J_{\boldsymbol\theta}(\z) \tilde{f}^* ( \boldsymbol\theta(\z))$ a.e. on $\partial\Omega.$
% % Since $\tilde{ f}^* \in L^\infty(\partial\boldsymbol \theta(\Omega))$ and $J_{\boldsymbol \theta}$ is a polynomial, $J_{\boldsymbol\theta} \tilde{f}^* \circ \boldsymbol\theta \in L^\infty(\partial\Omega)$ and anti-symmetric with respect to the group $G.$ 
% % Thus, $J_{\boldsymbol\theta} \tilde{f}^* \circ \boldsymbol\theta \in H^\infty(\partial\Omega)$ and consequently, $J_{\boldsymbol\theta} f \circ \boldsymbol\theta \in H^\infty(\Omega)$ and $$\|J_{\boldsymbol\theta} f \circ \boldsymbol\theta\|_{\infty,\Omega}=\|J_{\boldsymbol\theta} \tilde{f}^* \circ \boldsymbol\theta\|_{L^\infty(\Omega)}.$$ 

% Since $g$ and $J_{\boldsymbol \theta}$ are bounded and holomorphic on $\Omega$, $\frac{g}{J_{\boldsymbol \theta}}$ is a well defined holomorphic function on $\Omega \setminus N(J_{\boldsymbol \theta})$, where $N(J_{\boldsymbol \theta})$ is the zero set of $J_{\boldsymbol \theta}.$ 
% By the Riemann removable singularity theorem \cite[Theorem 4.2.9]{JP2000}, one can extend $\frac{g}{J_{\boldsymbol \theta}}$ as a holomorphic function on $\Omega$ which we denote by $F$ that satisfies the condition $g=J_{\boldsymbol \theta} F.$
% Furthermore, we can show that for almost every $\z \in \partial\Omega$, $\lim_{r \to 1^-} F(r\z)=\tilde{f}^*(\boldsymbol \theta(\z))$ using the fact that $\lim_{r\to 1^-} g(r\z)=g^*(\z)$ almost everywhere on $\partial\Omega.$ \textcolor{blue}{check}.

% Now, if we can prove that $F$ is a bounded holomorphic function on $\Omega$ then $F$ being $G$-invariant, $F$ must factor through a bounded holomorphic function $f$ on $\boldsymbol \theta(\Omega)$ such that $F=f\circ \boldsymbol \theta.$ It can be easily verified that $i'(f)=\tilde{f}^*$ and $\|f\|_{\infty, \boldsymbol \theta(\Omega)}=\|\tilde{f}^*\|_{L^\infty(\partial\boldsymbol \theta(\Omega))}.$
% condition
% Thus, it remains to show that $F$ is bounded on $\Omega.$ To this begin, we first prove that $F$ is in Nevanlinna class $N(\Omega)$, that is, $$\sup_{0<r<1}\int_{\partial\Omega} \log^+ |F(r\z)| d\mu(\z) <\infty.$$
% Note that as $G$ and $J_{\boldsymbol \theta}$ is bounded and holomorphic on $\Omega$, we get $$\log^+ \left|\frac{G}{J_{\boldsymbol \theta}}\right| \leq \log^+|G|+\log^-|J_{\boldsymbol \theta}| <\infty.$$
% Thus, the family $\{ \log^+|f_r|: 0<r<1\}$ is uniformly integrable. In fact, it satisfies 
% $$\lim_{r \to 1^-} \int_{\partial\Omega} \log^+|F(r\z)|d\mu(\z)=\int_{\partial\Omega} \log^+|\tilde{f}^*(\boldsymbol{\theta}(\z))d\mu(\z).$$
% Thus, by \cite[Lemma 1]{SM1976}, $$\log|F(\z)|\leq \int_{\partial\Omega} P(\z,\boldsymbol\zeta) \log|\tilde{f}^* (\boldsymbol \theta(\boldsymbol \zeta))| d\mu(\boldsymbol \zeta).$$
% Recall that $\tilde{f}^*\circ \boldsymbol \theta$ is essentially bounded on $\partial\Omega$ w. r. to the measure $\mu$. If $\|\tilde{f}^* \circ \boldsymbol \theta\|_{L^\infty (\partial\Omega)} \leq M,$ then 
% $$\log|F(\z)| \leq \log M \int_{\partial\Omega}P(\z, \boldsymbol \zeta) d\mu(\boldsymbol \zeta)=\log M.$$ Here, $P(\z, \boldsymbol \zeta)$ is the Poisson kernel of the domain $\Omega.$
% Thus, $F$ is bounded.
    
% Since the zero set of $J_{\boldsymbol\theta}$ has Lebesgue measure zero on $ \partial\Omega$, we deduce that $\tilde{f}^*(\boldsymbol\theta(\z))=\frac{g^*(\z)}{J_{\boldsymbol\theta}(\z)}$ a.e on $\partial\Omega.$
% Additionally, we note that for almost every $\z \in \mathbb T^d$,  $$g^*(\z)=\lim_{r\to 1^-} J_{\\xi}(r\z) f\circ\\xi (r\z)=J_{\\xi}(\z) \lim_{r \to 1^-} f(\\xi(r\z))= J_{\\xi}(\z) (f\circ \\xi)^*(\z).$$ This implies that $\|(f\circ \\xi)^*\|_{L^\infty(\mathbb T^d)}=\|\tilde{f}^* \circ \boldsymbol \xi\|_{L^\infty(\mathbb T^d)}=\|\tilde{f}^*\|_{L^\infty(\partial\mathbb G_d)}.$
% Because $f$ is holomorphic on $\mathbb G_d$, the composition $f\circ \\xi$ is holomorphic on $\mathbb D^d$. 
% Since \((f\circ\boldsymbol\xi)^*\in L^\infty(\mathbb T^d)\), the Poisson integral formula implies that \(f\circ\boldsymbol\xi\in H^\infty(\mathbb D^d)\). Hence \(f\in H^\infty(\mathbb G_d)\) and satisfies
% $$ \|f\|_{\infty, \mathbb G_d}= \|f\circ \boldsymbol \xi\|_{\infty, \mathbb D^d}= \|(f\circ \boldsymbol\xi)^*\|_{L^\infty(\mathbb T^d)}= \|\tilde{f}^* \|_{L^\infty(\partial\mathbb G_d)}.$$
% This proves the asserted norm equality and shows that $i'(H^\infty(\mathbb G_d))=H^\infty(\partial\mathbb G_d).$






\section{Commutant lifting}\label{S3}
In this section, we characterize contractive module maps on quotient modules of \(H^2(\boldsymbol \theta(\Omega))\) that admit Schur-class lifts. 
More precisely, given a quotient module $\mathcal Q$ of $H^2(\boldsymbol \theta(\Omega))$, we classify all contractive module maps on $\mathcal Q$ that arise as compressions $S_{\varphi}=P_{\mathcal{Q}}T_{\varphi}|_{\mathcal{Q}}$ for some $\varphi\in \mathcal S(\boldsymbol \theta(\Omega))$.

By Proposition \ref{required result 1}, the Hardy space $H^2(\boldsymbol \theta(\Omega))$ is canonically identified with its boundary realization $H^2(\partial\boldsymbol \theta(\Omega)).$ 
%via the boundary value map $f \longmapsto \tilde f^{*}.$
Accordingly, we make no distinction between a function in $H^2(\boldsymbol \theta(\Omega))$ and its boundary realization unless necessary. Thus, every quotient module of $H^2(\boldsymbol \theta(\Omega))$ may be regarded as a closed subspace of $H^2(\partial\boldsymbol \theta(\Omega))$.
Henceforth, $\mathcal Q$ denotes a nonzero quotient module of $H^2(\partial\boldsymbol \theta(\Omega))$. 
Since $\mathcal{Q}^\perp$ is a submodule, the nontriviality of $\mathcal{Q}$ implies that $P_{\mathcal Q}1\neq 0.$
Consequently, \beq \label{Q description}\mathcal Q = \overline{\operatorname{span}}\left\{P_{\mathcal Q}\boldsymbol \zeta^\alpha P_{\mathcal Q}1: \alpha\in \mathbb Z_+^d\right\}.\eeq
For a closed subspace $\mathcal E\subseteq L^2(\partial\boldsymbol \theta(\Omega))$, set
$\mathcal E^{conj}=\{\bar{f}:\, f\in \mathcal E\}.$
Further, let $$H_0^2(\partial\boldsymbol \theta(\Omega))=\left\{f\in H^2(\partial\boldsymbol \theta(\Omega)):\langle f,1\rangle_{H^2(\partial\boldsymbol \theta(\Omega))}=0\right\}.$$
Under the boundary identification, this corresponds to the subspace of functions in $H^2(\boldsymbol \theta(\Omega))$ vanishing at the origin.
% Since $\mathbb S_{\boldsymbol \theta(\Omega)}(\cdot,0)=1$, we have $f(0) =\langle f,\mathbb S_{\boldsymbol \theta(\Omega)}(\cdot,0)\rangle=\langle \tilde f^*,1\rangle_{H^2(\partial\boldsymbol \theta(\Omega))}$ for every $f\in H^2(\boldsymbol \theta(\Omega)).$
% Hence, $$i'^{-1}\big(H_0^2(\partial\boldsymbol \theta(\Omega))\big)=\{f\in H^2(\boldsymbol \theta(\Omega)): f(0)=0\}.$$
%We shall use the following elementary orthogonality relation.
The following orthogonality relation will be used repeatedly.
\begin{lemma}\label{orthogonality}
    $H^2_0(\partial\boldsymbol \theta(\Omega))\perp H^2(\partial\boldsymbol \theta(\Omega))^{conj}.$
\end{lemma}
\begin{proof}
Let $\tilde f^*\in H_0^2(\partial\boldsymbol \theta(\Omega))$ and $\tilde g^*\in H^2(\partial\boldsymbol \theta(\Omega))$, and let $f\in H^2_0(\boldsymbol \theta(\Omega))$ and $g\in H^2(\boldsymbol \theta(\Omega))$ denote the associated interior functions. Since $\tilde f_r\tilde g_r \longrightarrow \tilde f^*\tilde g^*$ in $L^1(\partial\boldsymbol \theta(\Omega))$ as $r\to 1^-,$
it follows that \beq \label{eq 6}\langle \tilde f^*,\overline{\tilde g^*}\rangle_{L^2(\partial\boldsymbol \theta(\Omega))} = \lim_{r\to1^-} \frac{1}{c_{\boldsymbol \theta}^2}\int_{\partial\boldsymbol \theta(\Omega)} \tilde f_r\tilde g_r\, d\mu_{\boldsymbol \theta}=\frac{1}{c_{\boldsymbol \theta}^2}\lim_{r\to 1^-} \int_{\partial\Omega} |J_{\boldsymbol \theta}|^2 f(\boldsymbol \theta(r\z))g(\boldsymbol \theta(r\z)) d\mu(z).\eeq 
For a fixed $0<r<1$, define $H(\z)=f(\boldsymbol \theta(r\z))g(\boldsymbol \theta(r\z))$. Since the map $ \Omega\ni \z \mapsto H(\z)$ is holomorphic in a neighborhood of $\overline{\Omega},$ it admits a homogeneous expansion
 $H(\z)=\sum_{j=0}^\infty H_j(\z)$, where $H_0(\z)=f(\boldsymbol \theta(0)) g(\boldsymbol \theta(0)).$ Then $$\int_{\partial\Omega} H(\z)|J_{\boldsymbol \theta}|^2 d\mu(\z)= \sum_{j=0}^\infty \int_{\partial\Omega} H_j(\z) |J_{\boldsymbol \theta}(\z)|^2 d\mu(\z).$$
Since $\mu$ is rotation invariant and $J_{\boldsymbol \theta}$ is a homogeneous polynomial, $|J_{\boldsymbol \theta}(e^{it} \z)|=|J_{\boldsymbol \theta}(\z)|$ and $H_j(e^{it}\z)=e^{ijt} H_{j}(\z).$ 
% Thus, $$\int_{\partial \Omega} H_{j}(\z)|J_{\boldsymbol \theta}(\z)|^2 d\mu(\z)=\int_{\partial \Omega} H_{j}(e^{it}\z)|J_{\boldsymbol \theta}(e^{it}\z)|^2 d\mu(\z)=e^{ijt}\int_{\partial \Omega} H_{j}(\z)|J_{\boldsymbol \theta}(\z)|^2 d\mu(\z).$$
%Since this holds for any real $t$, 
Thus, we have $$\int_{\partial \Omega} H_{j}(\z)|J_{\boldsymbol \theta}(\z)|^2 d\mu(\z)=\begin{cases}
    c_{\boldsymbol\theta}^2 f(\boldsymbol\theta(0))g(\boldsymbol\theta(0)), \quad j=0\\
    0, \qquad\qquad\qquad\;\;\;\;\;\; j\neq 0
\end{cases}.$$
Consequently, $$\int_{\partial\Omega} H(\z)|J_{\boldsymbol \theta}|^2 d\mu(\z)=c_{\boldsymbol \theta}^2 f(\boldsymbol \theta(0)) g(\boldsymbol \theta(0)).$$
As $f(0)=0$ and $\boldsymbol \theta(0)=0$, \eqref{eq 6} yields $\langle \tilde f^*,\overline{\tilde g^*}\rangle=0.$
\end{proof}

It follows from Lemma \ref{orthogonality} that 
 $$ H^2(\partial\boldsymbol\theta(\Omega)) + H^2(\partial\boldsymbol\theta(\Omega))^{conj} = H_0^2(\partial\boldsymbol\theta(\Omega)) \oplus H^2(\partial\boldsymbol\theta(\Omega))^{conj}$$ is a closed subspace of $L^2(\partial\boldsymbol\theta(\Omega))$. 
Define $$\mathcal M_{\boldsymbol\theta(\Omega)} = L^2(\partial\boldsymbol\theta(\Omega)) \ominus \big( H^2(\partial\boldsymbol\theta(\Omega)) + H^2(\partial\boldsymbol\theta(\Omega))^{conj} \big).$$
By construction, the spaces $H^2_0(\partial\boldsymbol\theta(\Omega)),\;\mathcal M_{\boldsymbol\theta(\Omega)},$ and $\mathcal{Q}^{conj}$ are mutually orthogonal in $L^2(\partial\boldsymbol\theta(\Omega)).$ 
%making them disjoint in $L^1(\partial\boldsymbol\theta(\Omega)).$ 
Accordingly, we define their algebraic direct sum in $L^1(\partial\boldsymbol\theta(\Omega))$: 
\beq \label{eq 14} \mathcal M_{\mathcal Q} =\mathcal Q^{conj} \dotplus \mathcal M_{\boldsymbol\theta(\Omega)} \dotplus H_0^2(\partial\boldsymbol\theta(\Omega)).\eeq 
We now establish our first characterization of commutant lifting on quotient modules of $H^2(\boldsymbol\theta(\Omega))$ (c.f \cite[Theorem 4.1]{DeS2026}).
\begin{theorem}\label{1st characterization}
    Let $\mathcal{Q}\subseteq  H^2(\partial\boldsymbol\theta(\Omega))$ be a quotient module and let $X\in \mathcal{B}_1(\mathcal{Q})$ be a contractive module map. Then $X$ admits a Schur-class lift if and only if the functional $X_{\mathcal{Q}}: \left(\mathcal{M}_{\mathcal{Q}}, \|\cdot\|_{L^1({\partial\boldsymbol\theta(\Omega)})}\right) \to \mathbb C$ defined by \beq \label{eq 11} X_{\mathcal{Q}}(f)=\int_{\partial\boldsymbol\theta(\Omega)}\psi f d\mu_{\boldsymbol\theta},\qquad f\in \mathcal{M}_{\mathcal{Q}}\eeq is contractive, where $\psi=XP_{\mathcal{Q}}1.$
\end{theorem}
\begin{proof}
    Assume first that $X\in \mathcal{B}_1(\mathcal{Q})$ admits a Schur-class lift. Then there exists $\varphi\in \mathcal{S}(\boldsymbol\theta(\Omega))$ such that $X=S_{\varphi}=P_{\mathcal{Q}} T_{\varphi}|_{\mathcal{Q}}.$
    Since $\varphi\in H^\infty(\boldsymbol\theta(\Omega))$ and $\mathcal{Q}^\perp$ is a submodule, it follows that $\varphi\mathcal{Q}^\perp \subseteq \mathcal{Q}^\perp.$ Hence, \beq \label{eq 7} P_{\mathcal{Q}}\varphi= P_{\mathcal{Q}}T_\varphi P_{\mathcal{Q}}1+ P_{\mathcal{Q}}T_\varphi P_{\mathcal{Q}^\perp}1=P_{\mathcal{Q}}T_{\varphi}P_{\mathcal{Q}}1=S_{\mathcal{\varphi}} P_{\mathcal{Q}}1=X P_{\mathcal{Q}}1=\psi.\eeq
Define the functional $\chi_{\varphi}: L^1(\partial\boldsymbol\theta(\Omega)) \to \mathbb C$ by 
$$\chi_{\varphi}(f)=\int_{\partial\boldsymbol\theta(\Omega)} \varphi f d\mu_{\boldsymbol\theta} = \inp{\varphi}{\bar{f}}_{L^2(\partial\boldsymbol\theta(\Omega))},\qquad f\in L^1(\partial\boldsymbol\theta(\Omega)).$$
Since $\varphi \in \mathcal{S}(\boldsymbol\theta(\Omega))$, $\|\chi_{\varphi}\| = \| \varphi\|_{L^\infty(\partial\boldsymbol\theta(\Omega))} \leq 1.$ We claim that $\chi_{\varphi}|_{\mathcal{M}_{\mathcal{Q}}}=X_{\mathcal{Q}}.$ Indeed, this follows by checking separately on the three summands of $\mathcal{M}_{\mathcal{Q}}$ in \eqref{eq 14}.
 
 If $\bar{f}\in \mathcal{Q}^{conj}$, then
\beqn X_{\mathcal{Q}}(\bar{f}) = \int_{\partial\boldsymbol\theta(\Omega)} \psi \bar{f} d\mu_{\boldsymbol \theta} = \inp{\psi}{f}_{H^2(\partial\boldsymbol\theta(\Omega))} \overset{\eqref{eq 7}}{=} \inp{P_{\mathcal{Q}} \varphi}{f}_{H^2(\partial\boldsymbol\theta(\Omega))}
= \inp{\varphi}{f}_{L^2(\partial\boldsymbol\theta(\Omega))} = \chi_{\varphi}(\bar{f}).\eeqn
Note that $\mathcal{M}_{\boldsymbol\theta(\Omega)} = \mathcal{M}_{\boldsymbol\theta(\Omega)}^{conj}$, so if $f\in \mathcal{M}_{\boldsymbol\theta(\Omega)}$, its conjugate $\bar{f}$ is also in $\mathcal{M}_{\boldsymbol\theta(\Omega)}$. Since $\varphi, \psi \in H^2(\partial\boldsymbol\theta(\Omega))$ are orthogonal to $\mathcal{M}_{\boldsymbol\theta(\Omega)}$, 
$$ X_{\mathcal{Q}}(f) = \inp{\psi}{\bar{f}}_{L^2(\partial\boldsymbol\theta(\Omega))} = 0 = \inp{\varphi}{\bar{f}}_{L^2(\partial\boldsymbol\theta(\Omega))} = \chi_{\varphi}(f) \qquad \text{for all} \; f\in \mathcal{M}_{\boldsymbol\theta(\Omega)}. $$
Finally, if $f\in H^2_{0}(\partial\boldsymbol\theta(\Omega))$,
then Lemma \ref{orthogonality} yields 
$$ X_{\mathcal{Q}}(f)  = \overline{\inp{\overline{\psi}}{f}_{L^2(\partial\boldsymbol\theta(\Omega))}} = 0 = \overline{\inp{\overline{\varphi}}{f}_{L^2(\partial\boldsymbol\theta(\Omega))}} = \chi_{\varphi}(f). $$
Hence, $X_{\mathcal{Q}} = \chi_{\varphi}|_{\mathcal{M}_{\mathcal{Q}}}$, and therefore $\|X_{\mathcal{Q}}\| \leq \|\chi_{\varphi}\| \leq 1.$

Conversely, suppose $X_{\mathcal{Q}}$ is contractive on $\mathcal{M}_{\mathcal{Q}}.$ By the Hahn-Banach theorem, there exists a bounded linear functional $\widetilde{X}_{\mathcal{Q}}: L^1(\partial\boldsymbol\theta(\Omega))\to \mathbb C$ such that $\widetilde{X}_{\mathcal{Q}}|_{\mathcal{M}_{\mathcal{Q}}} = X_{\mathcal{Q}}$ and $\|\widetilde{X}_{\mathcal{Q}}\| = \|X_{\mathcal{Q}}\| \leq 1.$ 
By duality, %$(L^1(\partial\boldsymbol\theta(\Omega)))^*=L^\infty(\partial\boldsymbol\theta(\Omega))$, 
there exists $\varphi \in L^\infty(\partial\boldsymbol\theta(\Omega))$ such that $$\widetilde{X}_{\mathcal{Q}}(f) = \int_{\partial\boldsymbol\theta(\Omega)} \varphi f d\mu_{\boldsymbol \theta}=\inp{\varphi}{\bar{f}}_{L^2(\partial\boldsymbol\theta(\Omega))}, \qquad f \in L^1(\partial\boldsymbol\theta(\Omega)),$$
with $\|\varphi\|_{L^\infty(\partial\boldsymbol\theta(\Omega))} = \|\widetilde{X}_{\mathcal{Q}}\| \leq 1.$

We claim that $\varphi\in H^2(\partial\boldsymbol\theta(\Omega)).$ Since $$L^2(\partial\boldsymbol\theta(\Omega)) = H^2(\partial\boldsymbol\theta(\Omega)) \oplus H^2_0(\partial\boldsymbol\theta(\Omega))^{conj} \oplus \mathcal{M}_{\boldsymbol\theta(\Omega)},$$ it suffices to show $\varphi$ is orthogonal to the last two summands. If $f\in \mathcal{M}_{\boldsymbol\theta(\Omega)},$ then 
$$ \inp{\varphi}{f}_{L^2(\partial\boldsymbol\theta(\Omega))} = \widetilde{X}_{\mathcal{Q}}(\bar{f}) = X_{\mathcal{Q}}(\bar{f}) =  \inp{\psi}{f}_{L^2(\partial\boldsymbol\theta(\Omega))} = 0. $$
Similarly, if $\bar{h}\in H^2_0(\partial\boldsymbol\theta(\Omega))^{conj},$ then by Lemma \ref{orthogonality}
$$ \inp{\varphi}{\bar{h}}_{L^2(\partial\boldsymbol\theta(\Omega))} = \widetilde{X}_{\mathcal{Q}}(h) = X_{\mathcal{Q}}(h) = \overline{\inp{\overline{\psi}}{h}_{L^2(\partial\boldsymbol\theta(\Omega))}} = 0. $$
 Hence, $\varphi \in H^\infty(\partial\boldsymbol\theta(\Omega)).$ Since $\|\varphi\|_{L^\infty(\partial\boldsymbol\theta(\Omega))}\leq 1$, Proposition~\ref{required result 2} shows that $\varphi$ is the boundary realization of a Schur function on $\boldsymbol\theta(\Omega),$ which we continue to denote by $\varphi.$ 

Finally, we show that $X=S_{\varphi}.$ For any $f\in \mathcal{Q}$,
$$ \inp{\varphi}{f}_{H^2(\partial\boldsymbol\theta(\Omega))} = \widetilde{X}_{\mathcal{Q}}(\bar{f}) = X_{\mathcal{Q}}(\bar{f})= \inp{\psi}{f}_{H^2(\partial\boldsymbol\theta(\Omega))}. $$
Hence $P_{\mathcal{Q}} \varphi = \psi=XP_{\mathcal{Q}}1.$
Since $\mathcal{Q}^\perp$ is a submodule,
$$ XP_{\mathcal{Q}}1= P_{\mathcal{Q}}\varphi = P_{\mathcal{Q}} T_{\varphi}(P_{\mathcal{Q}}1 + P_{\mathcal{Q}^\perp}1) = P_{\mathcal{Q}} T_{\varphi}P_{\mathcal{Q}}1 = S_{\varphi} P_{\mathcal{Q}}1. $$
As $X$ commutes with compressed coordinate multipliers, $XS_{ \boldsymbol \zeta}^\alpha=S_{ \boldsymbol \zeta}^\alpha X$ for all $\alpha\in \mathbb Z_+^d.$ 
%where $S_{\boldsymbol \zeta}^\alpha= S_{\zeta_1}^{\alpha_1}\cdots S_{\zeta_d}^{\alpha_d}.$
Therefore, 
\beqn
    S_{\varphi}(P_{\mathcal{Q}} \boldsymbol \zeta^\alpha P_{\mathcal{Q}}1) = S_{\boldsymbol \zeta}^\alpha S_{\varphi} P_{\mathcal{Q}}1 
    &= S_{\boldsymbol \zeta}^\alpha X P_{\mathcal{Q}}1 = X S_{\boldsymbol \zeta}^\alpha P_{\mathcal{Q}}1 = X (P_{\mathcal{Q}} \boldsymbol \zeta^\alpha P_{\mathcal{Q}}1).
\eeqn
In view of \eqref{Q description}, we conclude that $X=S_{\varphi}.$ This completes the proof of the theorem.
\end{proof}


\begin{remark}
The contractivity assumption on $X$ in Theorem \ref{1st characterization} can be relaxed to simple boundedness. More precisely, if $X\in \mathcal{B}(\mathcal{Q})$ is a bounded module map such that the associated functional $X_{\mathcal{Q}}$ on $\mathcal{M}_{\mathcal{Q}}$, defined by \eqref{eq 11}, is contractive, then $X$ admits a Schur-class lift. In particular, $X$ is then necessarily contractive.
\end{remark}

Building on Theorem~\ref{1st characterization}, we next reformulate the contractivity of the functional $X_{\mathcal{Q}}$ as a distance condition in $L^1(\partial\boldsymbol\theta(\Omega)).$ 
Recall that $X_{\mathcal{Q}}$ is completely determined by the vector $\psi=XP_{\mathcal{Q}}1 \in \mathcal{Q}.$ It follows from the proof of Theorem \ref{1st characterization} that $X_{\mathcal{Q}}$ vanishes on $\mathcal{M}_{\boldsymbol\theta(\Omega)}\dotplus H^2_0(\partial\boldsymbol\theta(\Omega))$ and on the orthogonal complement of $\overline{\psi}$ in $\mathcal{Q}^{conj}.$
Hence, the kernel of $X_{\mathcal{Q}}$ is given by
$$\widetilde{\mathcal{M}}_{\mathcal{Q}, X}:=\mathrm{ker}X_{\mathcal{Q}}=(\mathcal{Q}^{conj}\ominus \mathrm{span}\{\overline{\psi}\}) \dotplus \mathcal{M}_{\boldsymbol\theta(\Omega)}\dotplus H^2_{0}(\partial\boldsymbol\theta(\Omega)).$$ 
This leads to the following geometric characterization of commutant lifting. 
%The following theorem characterizes commutant lifting in terms of the distance from the normalized vector $\overline{\psi}/\|\psi\|^2_{L^2(\partial\boldsymbol\theta(\Omega))}$ to $\widetilde{\mathcal M}_{\mathcal Q,X}$.
\begin{theorem}\label{2nd characterization}
    Let $\mathcal{Q}\subseteq H^2(\partial\boldsymbol\theta(\Omega))$ be a quotient module and let $X\in \mathcal{B}_1(\mathcal{Q})$ be a contractive module map. Set
$\psi=XP_{\mathcal Q}1.$ Then
    $X$ admits a Schur-class lift if and only if \beq \label{eq 15} \dist_{L^1(\partial\boldsymbol\theta(\Omega))}\left( \frac{\overline{\psi}}{\|\psi\|^2_{L^2(\partial\boldsymbol\theta(\Omega))}}, \widetilde{\mathcal{M}}_{\mathcal{Q}, X}\right) \geq 1.\eeq
\end{theorem}
\begin{proof}
    By the definition of the distance to a subspace, 
    $$ \dist_{L^1(\partial\boldsymbol\theta(\Omega))}\left( \frac{\overline{\psi}}{\|\psi\|^2_{L^2(\partial\boldsymbol\theta(\Omega))}}, \widetilde{\mathcal{M}}_{\mathcal{Q}, X}\right)=\inf \left\{\left\|\frac{\overline{\psi}}{\|\psi\|^2_{L^2(\partial\boldsymbol\theta(\Omega))}}+h\right\|_{L^1(\partial\boldsymbol\theta(\Omega))}: h\in \widetilde{\mathcal{M}}_{\mathcal{Q}, X} \right\}.$$ 
    Consequently, the distance condition in \eqref{eq 15} is equivalent to
   $$\|\lambda\overline{\psi}+h\|_{L^1(\partial\boldsymbol\theta(\Omega))} \geq |\lambda|\|\psi\|^2_{L^2(\partial\boldsymbol\theta(\Omega))}$$ for every $h \in \widetilde{\mathcal{M}}_{\mathcal{Q}, X}$ and every nonzero $\lambda \in \mathbb C.$
   
   Because $\mathcal{M}_{\mathcal{Q}} = \mathbb{C}\overline{\psi} \dotplus \widetilde{\mathcal{M}}_{\mathcal{Q}, X}$, every $f\in \mathcal{M}_{\mathcal Q}$ decomposes uniquely as $f = \lambda \overline{\psi} + h$ for some $\lambda\in\mathbb C$ and $h\in \widetilde{\mathcal{M}}_{\mathcal Q,X}$. Since $\mathrm{ker} X_{\mathcal{Q}}=\widetilde{\mathcal{M}}_{\mathcal Q,X}$,
\beq \label{eq 9} X_{\mathcal{Q}}(\lambda \overline{\psi} + h) = \lambda \int_{\partial\boldsymbol\theta(\Omega)} \psi \overline{\psi} d\mu_{\boldsymbol \theta} + X_{\mathcal{Q}}(h) = \lambda \|\psi\|^2_{L^2(\partial\boldsymbol\theta(\Omega))}. \eeq


If $X$ admits a Schur-class lift, Theorem \ref{1st characterization} gives that $X_{\mathcal{Q}}$ is contractive on $\mathcal{M}_{\mathcal{Q}}$. Thus, for every $\lambda \in \mathbb{C}$ and $h \in \widetilde{\mathcal{M}}_{\mathcal{Q}, X}$,
$$\|\lambda\overline{\psi} + h\|_{L^1(\partial\boldsymbol{\theta}(\Omega))} \geq |X_{\mathcal{Q}}(\lambda\overline{\psi} + h)| \overset{\eqref{eq 9}}{=} |\lambda| \|\psi\|^2_{L^2(\partial\boldsymbol{\theta}(\Omega))}.$$

% Assume first that $X$ admits a Schur-class lift. By Theorem~\ref{1st characterization}, the functional $X_{\mathcal Q}$ is contractive on $\mathcal{M}_{\mathcal Q}.$ Hence, for every $\lambda \in \mathbb C$ and every $h\in \widetilde{\mathcal{M}}_{\mathcal Q,X}$,
% $$ \|\lambda\overline{\psi} + h\|_{L^1(\partial\boldsymbol\theta(\Omega))} \geq |X_{\mathcal{Q}}(\lambda\overline{\psi} + h)| \overset{\eqref{eq 9}}{=} |\lambda|\|\psi\|^2_{L^2(\partial\boldsymbol\theta(\Omega))}. $$

Conversely, suppose that $\|\lambda\overline{\psi}+h\|_{L^1(\partial\boldsymbol\theta(\Omega))} \geq |\lambda|\|\psi\|^2_{L^2(\partial\boldsymbol\theta(\Omega))}$ for every nonzero $\lambda \in \mathbb C$ and $h\in \widetilde{\mathcal{M}}_{\mathcal{Q}, X}.$
Let $f = \lambda \overline{\psi} + h$ be an arbitrary element of $\mathcal{M}_{\mathcal{Q}}$. If $\lambda=0$, then $X_{\mathcal{Q}}(f) = 0 \leq \|f\|_{L^1(\partial\boldsymbol\theta(\Omega))}$. If $\lambda \neq 0$, then
$$ \|f\|_{L^1(\partial\boldsymbol\theta(\Omega))} = \|\lambda \overline{\psi} + h\|_{L^1(\partial\boldsymbol\theta(\Omega))} \geq |\lambda| \|\psi\|^2_{L^2(\partial\boldsymbol\theta(\Omega))}\overset{\eqref{eq 9}}{=} \left|X_{\mathcal{Q}}(\lambda \overline{\psi}+h)\right|.$$
Thus $X_{\mathcal Q}$ is contractive on $\mathcal{M}_{\mathcal Q}$,  and the conclusion follows from Theorem~\ref{1st characterization}.
\end{proof}


\subsection{Polydisc}
It was shown by Rudin \cite{Rd1983} that Schur functions on the Euclidean unit ball $\mathbb B_d$ can be approximated, in the weak-$*$ topology, by inner functions. More recently, this approximation phenomenon was used to obtain a characterization of commutant lifting on $\mathbb B_d$ in terms of inner functions (see \cite{BDS2025}). Motivated by these developments, we establish an analogous characterization of Schur-class liftability for quotient domains arising from the polydisc. We begin by fixing the relevant class of quotient domains and by formulating the corresponding notion of inner function.


Throughout this subsection, we take $\Omega=\mathbb D^d$, and let $G=G(m,n,d)$ be one of the imprimitive finite complex reflection groups in the Shephard-Todd classification \cite{ST1954}, where $m$ and $n$ are positive integers such that $n$ divides $m$. Let $\alpha$ be a primitive $m$-th  root of unity. The group $G(m,n,d)$ consists of transformations of the form 
$$\mathbb C^d \ni \z=(z_1,\ldots,z_d)\longmapsto (\alpha^{k_1} z_{\tau(1)}, \ldots, \alpha^{k_d} z_{\tau(d)}),$$ where $\tau\in \Sigma_d$ and $k_1,\ldots, k_d$ are integers such that $k_1+\cdots+k_d$ is divisible by $n.$

The associated basic polynomial map $\boldsymbol\theta=(\theta_1,\ldots,\theta_d):\mathbb D^d\to \boldsymbol \theta(\mathbb D^d)$ can be given explicitly: for each $1\leq j\leq d-1$, $\theta_j$ is the $j$-th elementary symmetric polynomial in $z_1^m,\ldots,z_d^m$, and $\theta_d(\boldsymbol z)=(z_1\cdots z_d)^{m/n}.$ 
The map $\boldsymbol\theta:\mathbb D^d\longrightarrow \boldsymbol\theta(\mathbb D^d)$ is proper and is factored by $G(m,n,d)$. In the special case $m=n=1$, the quotient domain $\boldsymbol\theta(\mathbb D^d)$ is the symmetrized polydisc. For notational simplicity, we shall write \(G\) for \(G(m,n,d)\) throughout the rest of this subsection.


Recall that the Shilov boundary of the quotient domain $\boldsymbol\theta(\mathbb D^d)$ is $\partial\boldsymbol \theta(\mathbb D^d)=\boldsymbol \theta(\mathbb T^d).$ In analogy with the polydisc, we say that a function $f\in H^\infty(\boldsymbol \theta(\mathbb D^d))$ is inner if the corresponding boundary function $\tilde{f}^*\in H^\infty(\boldsymbol  \theta(\mathbb T^d))$ is unimodular $\mu_{\boldsymbol \theta}$-a.e. on $\boldsymbol \theta(\mathbb T^d).$ 
%almost everywhere with respect to the measure $\mu_{\boldsymbol \theta}.$ 
Let $$\mathcal{I}(\boldsymbol \theta(\mathbb D^d))=\{f\in H^\infty(\boldsymbol \theta(\mathbb D^d)): f\circ\boldsymbol \theta \text{ is an inner function on } \mathbb D^d\}.$$ 
By Proposition \ref{required result 2}, for every $f\in H^\infty(\boldsymbol \theta(\mathbb D^d))$, $\tilde{f}^* \circ \boldsymbol \theta=(f \circ \boldsymbol \theta)^*$ almost everywhere on $\mathbb T^d$ and therefore the class of inner functions on $\boldsymbol \theta(\mathbb D^d)$ coincides with $\mathcal{I}(\boldsymbol \theta(\mathbb D^d)).$

We first prove that every Schur function on $\boldsymbol \theta(\mathbb D^d)$ is a weak-$*$ limit, in $L^\infty(\boldsymbol \theta(\mathbb T^d)),$ of inner functions on $\boldsymbol \theta(\mathbb D^d).$ 
To achieve this, we establish a preliminary result analogous to \cite[Theorem 5.5.1]{Rd1969}. The proof follows essentially the same line of argument as in Rudin’s theorem, with suitable modifications to incorporate the $G$-invariance condition. Recall that a function $f$ is $G$-invariant if $\sigma(f)=f$ for all $\sigma \in G.$
\begin{proposition}\label{caratheory}
    Let $f \in \mathcal{S}(\mathbb D^d)$ be $G$-invariant. Then there exists a sequence of $G$-invariant inner functions in the polydisc algebra $\mathcal{A}(\mathbb D^d)$ converging to $f$ uniformly on compact subsets of $\mathbb D^d.$ 
\end{proposition}
\begin{proof}
Fix a compact subset $K\subseteq \mathbb D^d$ and let $0<\epsilon<\frac{1}{2}.$ To enforce $G$-invariance, we define the orbit $\tilde{K}=\cup_{\sigma\in G} \sigma(K)$, which remains compact.
Choose a polynomial $P'$ such that $|P'(\z)| <1$ on $\overline{\mathbb D}^d$ and $|f(\boldsymbol{z}) - P'(\boldsymbol{z})| < \epsilon$ for all $\boldsymbol{z} \in \tilde{K}$.
To restore $G$-invariance, we consider the polynomial
$$ P(\boldsymbol{z}) = \frac{1}{|G|} \sum_{\sigma \in G} P'(\sigma(\z)). $$ Then $P$ is $G$-invariant and satisfies $|P(\z)| <1$  on $\overline{\mathbb D}^d.$ Moreover, for any $\z \in \tilde{K}$, we have
\beq \label{eq 17} |f(\boldsymbol{z}) - P(\boldsymbol{z})| \leq \frac{1}{|G|} \sum_{\sigma \in G} |f(\boldsymbol{z}) - P'(\sigma(\boldsymbol{z}))| = \frac{1}{|G|} \sum_{\sigma \in G} |f(\sigma(\boldsymbol{z}) - P'(\sigma(\boldsymbol{z})| < \epsilon, \eeq 
where we used the fact that $f$ is $G$-invariant and $\sigma(\z)\in \tilde{K}.$ For any polynomial $p$ in $d$ complex variables, let $\tilde{p}$ denote the polynomial defined by $\tilde{p}(\z)=\overline{p(\bar{\z})}.$  
Then, arguing as in \cite[Theorem~5.5.1]{Rd1969}, we obtain a monomial $M(\z)$ of the form $e^{i\nu} (z_1z_2\cdots z_d)^l$ of sufficiently high degree such that $m$ divides $l$ and satisfies $(i)$ $|M(\z)|=1$ on $\mathbb T^d$, $(ii)$ $M(\z) \tilde{P}(\frac{1}{\z})$ is a polynomial, and $(iii)$ both $|P(\z)M(\z)\tilde{P}(\frac{1}{\z})| <\epsilon$ and $|M(\z)|<\epsilon$ on $\tilde{K}.$ 
Define the rational function $$g(\z)=\frac{P(\z)+M(\z)}{1+M(\z)\tilde{P}(\frac{1}{\z})}.$$ Since $\alpha$ is $m$-th primitive root of unity and $l$ is divisible by $m$, then $\alpha^{l(k_1+\cdots+k_d)}=1$. Consequently, $M(\z)$ is $G$-invariant, that is, $M(\sigma(\z))=M(\z)$ for all $\sigma\in G.$ 
It is now straightforward to verify that $g$ is a $G$-invariant rational inner function in $\mathcal{A}(\mathbb D^d)$. Moreover, using \eqref{eq 17}, we obtain $|f(\z)-g(\z)| <5\epsilon$ for all $\z \in K.$ Since $K$ and $\epsilon$ were arbitrary, the result follows.
\end{proof}
 We now show that inner functions are weak-$*$ dense in the Schur class on $\boldsymbol \theta(\mathbb D^d).$

\begin{proposition}\label{weak*}
    Let  $\varphi\in \mathcal{S}(\boldsymbol \theta(\mathbb D^d)).$ Then there exists a sequence  $\{f_n\}\subseteq \mathcal{I}(\boldsymbol \theta(\mathbb D^d))$ converging to $\varphi$ in the weak-$*$ topology of $L^\infty(\boldsymbol \theta(\mathbb T^d)).$ 
\end{proposition}
\begin{proof}
    Let $\Phi=\varphi \circ \boldsymbol\theta.$ Then $\Phi$ is a $G$-invariant Schur function on $\mathbb D^d.$ By Proposition \ref{caratheory}, there exists a sequence $\{h_n\}$ of $G$-invariant inner functions in $\mathcal{A}(\mathbb D^d)$ such that $h_n \to \Phi$ uniformly on compact subsets of $\mathbb D^d.$ 
    By the analytic Chevalley-Shephard-Todd theorem (\cite[Theorem 3.1]{BDGR2022}), each $h_n$ factors through quotient map, that is, there exists $f_n\in H^\infty(\boldsymbol \theta(\mathbb D^d))$ such that $h_n =f_n \circ \boldsymbol \theta.$ 
    %for some bounded holomorphic function $f_n$ on $\boldsymbol \theta(\mathbb D^d).$
    Since $h_n$ is inner, $f_n \in \mathcal{I}(\boldsymbol \theta(\mathbb D^d)).$
   
    We claim that $f_n\to \varphi$ in the weak-$*$ topology of $L^\infty(\boldsymbol \theta(\mathbb T^d)).$
    By the change of variables formula, it suffices to show that
    $$ \int_{\mathbb T^d} h_n u d\mu \to \int_{\mathbb T^d} \Phi u d\mu \qquad \text{for all}\; u \in L^1(\mathbb T^d).$$
    % Since, for every $g\in L^1(\boldsymbol \theta(\mathbb T^d))$, we have $|J_{\boldsymbol \theta}|^2(g \circ \boldsymbol \theta) \in L^1(\mathbb T^d),$ it follows from the change of variables formula, together with the identities $h_n = f_n \circ \boldsymbol{\theta}$ and $\Phi = \varphi \circ \boldsymbol{\theta}$, that it suffices to prove $$ \int_{\mathbb T^d} h_n u d\mu \to \int_{\mathbb T^d} \Phi u d\mu \qquad \text{for all}\; u \in L^1(\mathbb T^d).$$
    Since trigonometric polynomials are dense in $L^1(\mathbb T^d)$ and the functions $h_n$ are uniformly bounded, it is enough to verify that

    \beq \label{eq 13}  \int_{\mathbb{T}^d} h_n(\boldsymbol{z}) \boldsymbol{z}^\alpha \bar{\boldsymbol{z}}^\beta \, d\mu(\boldsymbol{z}) \to \int_{\mathbb{T}^d} \Phi(\boldsymbol{z}) \boldsymbol{z}^\alpha \bar{\boldsymbol{z}}^\beta \, d\mu(\boldsymbol{z}) \qquad \text{for all}\; \alpha , \beta \in \mathbb Z_+^d.\eeq
    The integrals in \eqref{eq 13} are partial derivatives of $h_n$ and $\Phi$ at the origin. The convergence of these integrals follows directly from the Cauchy integral formula and the uniform convergence of $h_n \to \Phi$ on compact subsets of $\mathbb{D}^d$.
    \end{proof} 
    
    The weak-$*$ density of inner functions established above yields the following approximation characterization for commutant lifting. 
    For a subset $S\subseteq L^1(\boldsymbol \theta(\mathbb T^d))$ and a function $\psi \in H^2(\boldsymbol \theta(\mathbb D^d))$, we say a sequence $\{f_n\}\subseteq \mathcal{S}(\boldsymbol \theta(\mathbb D^d))$ converges to $\psi$ in the weak-$*$ topology determined by $S$ if
$$\lim_{n\to\infty}\int_{\boldsymbol \theta(\mathbb T^d)} f_n g\, d\mu_{\boldsymbol\theta}=\int_{\boldsymbol \theta(\mathbb T^d)} \psi g\, d\mu_{\boldsymbol\theta}\qquad \text{for all } g\in S.$$    

\begin{theorem}\label{3rd characterization}
    Let  $\mathcal{Q}\subseteq H^2(\boldsymbol \theta(\mathbb T^d))$ be a quotient module and let $X\in \mathcal{B}_1(\mathcal{Q})$ be a contractive module map. Then, $X$ admits a lift if and only if there exists a sequence $\{f_n\}\subseteq \mathcal{I}(\boldsymbol \theta(\mathbb D^d))$ converging to $\psi=XP_{\mathcal{Q}}1$ in the weak-$*$ topology determined by $\mathcal{M}_{\mathcal{Q}}.$
\end{theorem}
\begin{proof}
    Assume first that $X\in \mathcal{B}_1(\mathcal{Q})$ admits a Schur-class lift. By Theorem \ref{1st characterization}, there exists $\varphi \in \mathcal{S}(\boldsymbol \theta(\mathbb D^d))$ such that the associated functional $X_{\mathcal{Q}}$ on $\mathcal{M}_{\mathcal{Q}}$ is given by integration against $\varphi.$ Equivalently, for every $g\in \mathcal{M}_{\mathcal{Q}}$, $$\int_{ \boldsymbol \theta(\mathbb T^d)} \psi g d\mu_{\boldsymbol \theta}=X_{\mathcal{Q}} g=\int_{\boldsymbol \theta(\mathbb T^d)} \varphi g d\mu_{\boldsymbol \theta}.$$   
    By Proposition \ref{weak*}, there exists a sequence $\{f_n\}\subseteq \mathcal{I}(\boldsymbol \theta(\mathbb D^d))$ converging to $ \varphi$ in the weak-$*$ topology of $L^\infty(\boldsymbol \theta(\mathbb T^d)).$ In particular, for every $g\in \mathcal{M}_{\mathcal{Q}},$
    $$\lim_{n\to\infty}\int_{\boldsymbol \theta(\mathbb T^d)} f_n g\, d\mu_{\boldsymbol \theta}
    =\int_{\boldsymbol \theta(\mathbb T^d)} \varphi g\, d\mu_{\boldsymbol \theta}
    =\int_{\boldsymbol \theta(\mathbb T^d)} \psi g\, d\mu_{\boldsymbol \theta}.$$
    Thus, $f_n \to \psi$ in the weak-$*$ topology determined by $\mathcal{M}_{\mathcal{Q}}.$
    
    Conversely, suppose there exists a sequence $\{f_n\}\subseteq \mathcal{I}(\boldsymbol \theta(\mathbb D^d))$ such that
    $$\lim_{n\to\infty}\int_{\boldsymbol \theta(\mathbb T^d)} f_n g\, d\mu_{\boldsymbol \theta}
    =\int_{\boldsymbol \theta(\mathbb T^d)} \psi g\, d\mu_{\boldsymbol \theta}
    \qquad \text{for all } g\in \mathcal{M}_{\mathcal{Q}}.$$
    Since each $f_n$ is inner, $|f_n|=1$ $\mu_{\boldsymbol \theta}$-a.e. on $\boldsymbol \theta(\mathbb T^d)$.
    Thus, for every $g\in \mathcal{M}_{\mathcal{Q}},$
   \beqn 
        |X_{\mathcal{Q}}g|
        = \left|\int_{\boldsymbol \theta(\mathbb T^d)} \psi g\, d\mu_{\boldsymbol \theta}\right|= \left|\lim_{n\to\infty}\int_{\boldsymbol \theta(\mathbb T^d)} f_n g \, d\mu_{\boldsymbol \theta}\right|
        \leq\|g\|_{L^1(\boldsymbol \theta(\mathbb T^d))}. 
    \eeqn
    Hence $X_{\mathcal Q}$ is contractive on $\mathcal M_{\mathcal Q}.$ The conclusion now follows from Theorem \ref{1st characterization}.
    \end{proof}
    
The weak-$*$ convergence condition in Theorem~\ref{3rd characterization} can be further simplified. 
Recall from the proof of Theorem \ref{1st characterization} that the associated functional $X_{\mathcal{Q}}$ vanishes on $\mathcal{M}_{\boldsymbol \theta(\mathbb D^d)}\dotplus H^2_0(\boldsymbol \theta(\mathbb T^d)).$  
Consequently, convergence against $\mathcal{M}_{\mathcal{Q}}$  reduces exclusively to convergence determined by $\mathcal{Q}^{conj}.$ 
This observation leads to the following immediate corollary. 
\begin{corollary}\label{refined characterization}
    Let $\mathcal{Q} \subseteq H^2(\boldsymbol \theta(\mathbb T^d))$ be a quotient module and let $X\in \mathcal{B}_1(\mathcal{Q})$ be a contractive module map. Then, $X$ admits a lift if and only if there exists a sequence $\{f_n\}\subseteq \mathcal{I}(\boldsymbol \theta(\mathbb D^d))$ converging to $\psi=XP_{\mathcal{Q}}1$ in the weak-$*$ topology determined by $\mathcal{Q}^{conj}.$
\end{corollary}









%     Thus, the convergence of the integrals follows from the Cauchy integral formula on a polycircle together with the uniform convergence $h_n \to \Phi$ on compact subsets of $\mathbb D^d.$ 
% This completes the proof.


% Since $X_{\mathcal{Q}}$ vanishes on $\mathcal{M}_{\boldsymbol \theta(\mathbb D^d)}\dotplus H^2_0(\boldsymbol \theta(\mathbb T^d)),$ it is enough to test convergence against $\mathcal{Q}^{conj}.$


 % where the last inequality holds because $|f_n|=1$ almost everywhere on $\boldsymbol \theta(\mathbb T^d).$ 


% is the power series expansion 
% Recall that $J_{\boldsymbol \xi}(\z)=\prod_{1 \leq i <j \leq d} (z_j-z_i).$ Consequently, 
% $$|J_{\boldsymbol \xi}(\z)|^2= (d-1)! \sum_{i=1}^d |z_i|^2 +\sum_{\alpha \neq \beta }c_{\alpha, \beta}\z^\alpha \bar{\z}^\beta,$$
% %\substack{{\alpha, \beta\in\mathbb Z_+^d}\\ \alpha \neq \beta}
% where every term in the right-hand summation contains distinct holomorphic and antiholomorphic monomials. Using the orthogonality of distinct monomials in $L^2(\mathbb T^d),$ it follows that 
% $$\int_{\partial\mathbb G_d}
% \tilde f_r\tilde g_r\, d\mu_{\boldsymbol \xi}= \frac{1}{\sqrt{d!}} \int_{\mathbb T^d} |J_{\boldsymbol \xi}|^2 (\tilde{f}_r\circ\boldsymbol \xi) (\tilde{g}_r\circ\boldsymbol \xi) d\mu= \sqrt{d!}\int_{\mathbb T^d} f(\boldsymbol\xi(r\z)) g(\boldsymbol \xi(r\z)) d\mu(\z).$$
% Since $\z \mapsto f(\boldsymbol \xi(r\z))g(\boldsymbol \xi(r\z))$ is holomorphic on a neighborhood of $\overline{\mathbb D}^d$, the mean value property yields 
% $$\int_{\mathbb T^d} f(\boldsymbol \xi(r\z)) g(\boldsymbol \xi(r\z)) \, d\mu(\z) =f(\boldsymbol \xi(0))g(\boldsymbol \xi(0)).$$

%which is image of polydisc under a proper holomorphic map. 
%To this end, we first formalize the notion of inner functions on such quotient domains. 
%$\boldsymbol \theta(\mathbb D^d)$.


% \limsup_{n\to\infty}\int_{\partial\boldsymbol \theta(\mathbb D^d)} |f_n g|\, d\mu_{\boldsymbol \theta}= 

% the finite complex reflection group is $G(m,n,d)$ as appeared in the classifications of irreducible finite complex reflection groups by Shephard and Todd in \cite{ST1954}, where $m$ and $n$ are positive integers such that $n$ divides $m$. In fact, in this case, the proper holomorphic maps $\boldsymbol \theta=(\theta_1,\ldots. \theta_d)$ are as follows: for each $1\leq j \leq d-1$, $\theta_j$ are $j$-the elementary symmetric polynomial in $z_1^m, \ldots, z_d^m$ and $\theta_d(\z)=(z_1\cdots z_d)^{m/n}.$ 
% Furthermore, the group $ G(m,n,d)$ is generated by the complex reflections as follows: $\alpha$ is a $m$th primitive root of unity
% $$\mathbb C^d \ni \z=(z_1,\ldots,z_d)\mapsto (\alpha^{k_1} z_{\tau(1)}, \ldots, \alpha^{k_d} z_{\tau(d)})$$ for some $\tau \in \Sigma_d$ and integers $k_1,\ldots, k_d$ such that $k_1+\cdots+k_d$ is divisible by $n.$
% Note that for $m=1=n$, the domain $\boldsymbol \theta(\mathbb D^d)$ is nothing but the symmetrized polydisc.


\section{Interpolation} \label{S4}
In this section we apply the commutant lifting results of Section~\ref{S3} to interpolation problems on quotient domains. 
%$\boldsymbol \theta(\Omega)$ of a bounded symmetric domain $\Omega.$
This approach is motivated by Sarason’s operator-theoretic proof of the classical Nevanlinna–Pick interpolation theorem \cite[Section~4]{DS1967}, and by subsequent multivariable extensions for the polydisc and the Euclidean unit ball; see, for instance, \cite{DeS2026,BDS2025}.
Interpolation on the symmetrized bidisc, which is a proper image of the bidisc, was previously studied in \cite{AY2017,BS2018} using realization formulas for Schur functions. In contrast, our framework relies on commutant lifting and applies uniformly to quotient domains, including the symmetrized bidisc and the tetrablock.
%its subsequent extensions to multivariable settings such as the polydisc and the Euclidean unit ball (see, for instance, \cite{DeS2026,BDS2025}).
% Recall that the symmetrized bidisc $\mathbb G_2$ is a proper image of the bidisc $\mathbb D^2.$ Interpolation on $\mathbb G_2$ was previously investigated in \cite{AY2017,BS2018} using realization formulas for Schur functions. In contrast, our framework relies entirely on commutant lifting and leads to a different formulation of the interpolation problem that applies universally to quotient domains, including symmetrized bidisc and tetrablock.

The primary objective of this section is to characterize the existence of a function $\varphi\in \mathcal{S}(\boldsymbol \theta(\Omega))$ satisfying 
$$\varphi(\boldsymbol \zeta_i)=w_i, \qquad i=1,\ldots,p,$$
for a prescribed set of distinct points $\mathcal{Z}=\{\boldsymbol \zeta_1,\ldots,\boldsymbol \zeta_p\}\subseteq \boldsymbol \theta(\Omega)$ and target scalars $w_1,\ldots,w_p\in\mathbb D.$


We begin by recalling that $H^2(\boldsymbol \theta(\Omega))$ is a reproducing kernel Hilbert space with reproducing kernel $\mathbb S_{\boldsymbol \theta(\Omega)}$ is as given in \eqref{Hardy kernel}. For each $\zeta \in \boldsymbol \theta(\Omega)$, the kernel function $\mathbb S_{\boldsymbol \theta(\Omega)}(\cdot,\boldsymbol \zeta)$ satisfies
% with the reproducing kernel $\mathbb S_{\boldsymbol \theta(\Omega)}$ 
% given by \eqref{Hardy kernel}.
% For each $\boldsymbol \zeta \in \boldsymbol \theta(\Omega)$, the kernel function $\mathbb S_{\boldsymbol \theta(\Omega)}(\cdot,\boldsymbol \zeta): \boldsymbol \theta(\Omega) \to \mathbb C$ satisfies 
 the fundamental reproducing property
$$\inp{f}{\mathbb S_{\boldsymbol \theta(\Omega)}(\cdot,\boldsymbol \zeta)}_{H^2(\boldsymbol\theta(\Omega))}= f(\boldsymbol \zeta)\qquad \text{for all}\; f\in H^2(\boldsymbol \theta(\Omega)).$$
In particular, \beq \label{eq 10} T_{\zeta_i}^* \mathbb S_{\boldsymbol \theta(\Omega)}(\cdot, \boldsymbol \eta)=\bar{\eta}_i \mathbb S_{\boldsymbol\theta(\Omega)}(\cdot, \boldsymbol \eta), \qquad \boldsymbol \eta=(\eta_1,\ldots, \eta_d)\in \boldsymbol\theta(\Omega), \quad i=1,\ldots,d. \eeq
Given the set of $p$ distinct points $\mathcal{Z}$, we define the subspace
$$\mathcal{Q}_{\mathcal{Z}}=\mathrm{span}\{\mathbb S_{\boldsymbol\theta(\Omega)}(\cdot, \boldsymbol \zeta_i): i=1,\ldots, p\}.$$
It follows from \eqref{eq 10} that $\mathcal{Q}_{\mathcal{Z}}$ is a finite-dimensional quotient module of $H^2(\boldsymbol \theta(\Omega)).$ Its orthogonal complement is the space of functions vanishing on the interpolation nodes:
$$\mathcal{Q}_{\mathcal{Z}}^\perp =\left\{ g\in H^2(\boldsymbol \theta(\Omega)): g(\boldsymbol \zeta_i)=0 \text{ for all } i=1,\ldots,p\right\}.$$ We refer to quotient modules of this form as finite-dimensional {\it zero-based} quotient modules.
Since each kernel function $\mathbb S_{\boldsymbol \theta(\Omega)}(\cdot,\boldsymbol \zeta_j)$ is bounded, we have
 $\mathcal{Q}_{\mathcal{Z}} \subseteq H^\infty(\boldsymbol \theta(\Omega)).$

Let $X\in \mathcal{B}(\mathcal{Q}_{\mathcal{Z}})$ be a bounded module map,
and set $\psi =XP_{\mathcal{Q}_{\mathcal{Z}}} 1$. Then $ \psi \in \mathcal{Q}_{\mathcal{Z}}\subseteq H^\infty (\boldsymbol \theta(\Omega)).$
%Since $\mathcal{Q}_{\mathcal{Z}} =\overline{\mathrm{span}} \{P_{\mathcal{Q}_{\mathcal{Z}}} \boldsymbol \zeta^\alpha P_{\mathcal{Q}_{\mathcal{Z}}} 1: \alpha \in \mathbb Z_+^d \}$ and 
As $X$ is a module map, for any $\alpha \in \mathbb Z_+^d$ we obtain
\beqn X(P_{\mathcal{Q}_{\mathcal{Z}}} \boldsymbol \zeta^\alpha P_{\mathcal{Q}_{\mathcal{Z}}} 1) =X (S_{\boldsymbol \zeta}^\alpha  P_{\mathcal{Q}_{\mathcal{Z}}} 1)= S_{\boldsymbol \zeta}^\alpha X P_{\mathcal{Q}_{\mathcal{Z}}} 1= S_{\boldsymbol \zeta}^\alpha \psi=S_{\psi} P_{\mathcal{Q}_{\mathcal{Z}}} {\boldsymbol \zeta}^\alpha P_{\mathcal{Q}_{\mathcal{Z}}}1.\eeqn
Since $\mathcal{Q}_{\mathcal{Z}} =\overline{\mathrm{span}} \{P_{\mathcal{Q}_{\mathcal{Z}}} \boldsymbol \zeta^\alpha P_{\mathcal{Q}_{\mathcal{Z}}} 1: \alpha \in \mathbb Z_+^d \}$, it follows that $X=S_{\psi}$ on $\mathcal{Q}_{\mathcal{Z}}.$ 
%Thus every bounded module map on $\mathcal{Q}_{\mathcal{Z}}$ admits an $H^\infty(\boldsymbol \theta(\Omega))$-lift. The interpolation problem, however, requires such a lift to belong to the Schur class. The commutant lifting criteria from Section \ref{S3} therefore lead to the following interpolation theorem.
Thus, every bounded module map on $\mathcal{Q}_{\mathcal{Z}}$ admits an $H^\infty(\boldsymbol \theta(\Omega))$-lift.
%a lift in $H^\infty(\boldsymbol \theta(\Omega))$.
However, such a lift need not preserve the Schur-class norm structure on $\boldsymbol \theta(\Omega).$ 
%In this case, we say that X admits a {\it weak lift}.
Thus the interpolation problem reduces to deciding when the $H^\infty(\boldsymbol \theta(\Omega))$-lift can be chosen from the Schur class.
%To solve the interpolation problem, $X$ should admit a lift from the Schur class. 

Using the commutant lifting framework developed in Section~\ref{S3}, we now establish the following interpolation theorem.
%To solve the interpolation problem, it is necessary to pass from a weak lift to a contractive Schur-class lift. Using the commutant lifting framework developed in Section~\ref{S3}, we now establish the following interpolation theorem.
\begin{theorem}\label{interpolation 1st}
    Let $\mathcal{Z}=\{\boldsymbol \zeta_1,\ldots,\boldsymbol \zeta_p\}\subseteq \boldsymbol \theta(\Omega)$ be a set of $p$ distinct points, and let $w_1,\ldots,w_p\in \mathbb D.$
Define the scalars $c_1,\ldots,c_p$ via the Gram matrix inverse:
    $$\begin{bmatrix}
            c_1 \\  \vdots\\ c_p 
        \end{bmatrix}= {\begin{bmatrix}
            \mathbb S_{\boldsymbol \theta(\Omega)}(\boldsymbol \zeta_i, \boldsymbol \zeta_j)
        \end{bmatrix}}_{p\times p}^{-1} \begin{bmatrix}
            w_1 \\ \vdots \\w_p
        \end{bmatrix}.$$
    Set $\mathcal{M}_{\mathcal{Q}_{\mathcal{Z}}}=\mathcal{Q}_{\mathcal{Z}}^{conj} \dotplus \mathcal{M}_{\boldsymbol \theta(\Omega)}\dotplus H^2_{0}(\partial\boldsymbol \theta(\Omega))$ and $\psi_{\mathcal{Z}, \mathcal{W}}=\sum_{i=1}^p c_i \mathbb S_{\boldsymbol \theta(\Omega)}(\cdot, \boldsymbol \zeta_i). $ 
    Then there exists $\varphi\in \mathcal{S}(\boldsymbol \theta(\Omega))$ satisfying $\varphi(\boldsymbol \zeta_i)=w_i$ for all $i=1,\ldots,p$ if and only if the functional $X_{\mathcal{Q_{\mathcal{Z}}}}: \left(\mathcal{M}_{\mathcal{Q}_{\mathcal{Z}}}, \|\cdot \|_{L^1(\partial\boldsymbol \theta(\Omega))}\right) \to \mathbb C$ defined by $$X_{\mathcal{Q_{\mathcal{Z}}}} f=\int_{\partial\boldsymbol \theta(\Omega)} \psi_{\mathcal{Z}, \mathcal{W}} fd\mu_{\boldsymbol \theta}, \qquad f\in \mathcal{M}_{\mathcal{Q}_{\mathcal{Z}}}$$ is contractive.
     \end{theorem}
    %where $\psi_{\mathcal{Z}, \mathcal{W}}=\sum_{i=1}^p c_i \mathbb S_{\boldsymbol \theta(\Omega)}(\cdot, \boldsymbol \zeta_i). $
   
\begin{proof}
% Since the points $\boldsymbol\zeta_1,\ldots,\boldsymbol\zeta_p$ are distinct, the corresponding kernel functions are linearly independent, and hence the Gram matrix $G_{\mathcal Z}$ is invertible.
    Consider the operator $X_{\mathcal{Z,W}}\in \mathcal{B}(\mathcal{Q}_{\mathcal{Z}})$ defined by \beq \label{eq 12} X_{\mathcal{Z,W}}^* \mathbb S_{\boldsymbol \theta(\Omega)}(\cdot, \boldsymbol \zeta_i)=\overline{w}_i \mathbb S_{\boldsymbol \theta(\Omega)}(\cdot, \boldsymbol \zeta_i),\qquad i=1,\ldots,p.\eeq 
    A routine verification shows that $X_{\mathcal{Z},\mathcal{W}}$ is a bounded module map. Hence, by our preceding discussion, $\psi=X_{\mathcal{Z,W}}P_{\mathcal{Q}_{\mathcal{Z}}}1\in H^\infty(\boldsymbol \theta(\Omega))$ and $X_{\mathcal{Z,W}}=S_{\psi}$ on $\mathcal{Q}_{\mathcal{Z}}.$

    Suppose there exists $\varphi\in \mathcal{S}(\boldsymbol \theta(\Omega))$ such that $\varphi(\boldsymbol{\zeta}_i) = w_i$ for all $i = 1, \ldots, p$. Using the reproducing property,
    $$S_{\varphi}^* \mathbb S_{\boldsymbol \theta(\Omega)}(\cdot, \boldsymbol \zeta_i)=P_{\mathcal{Q}_{\mathcal{Z}}} T_{\varphi}^*|_{\mathcal{Q}_{\mathcal{Z}}} S_{\boldsymbol \theta(\Omega)}(\cdot, \boldsymbol \zeta_i) = \overline{\varphi(\boldsymbol \zeta_i)}\mathbb S_{\boldsymbol \theta(\Omega)}(\cdot, \boldsymbol \zeta_i)=\overline{w}_i \mathbb S_{\boldsymbol \theta(\Omega)}(\cdot, \boldsymbol \zeta_i).$$ 
    Comparing this with \eqref{eq 12}, we obtain $X_{\mathcal{Z,W}}=S_{\varphi}.$ Thus, $X_{\mathcal{Z},\mathcal{W}}$ admits a lift to the Schur function $\varphi.$
    
    Conversely, if $X_{\mathcal{Z,W}}=S_{\varphi}$ for some $\varphi\in \mathcal{S}(\boldsymbol \theta(\Omega)),$ then \eqref{eq 12} gives
    $$\overline{w}_i \mathbb S_{\boldsymbol \theta(\Omega)}(\cdot, \boldsymbol \zeta_i)=\overline{\varphi(\boldsymbol \zeta_i)} \mathbb S_{\boldsymbol \theta(\Omega)}(\cdot,\boldsymbol \zeta_i), \qquad i=1,\ldots,p,$$ and hence  $\varphi(\boldsymbol \zeta_i)=w_i$ for all $i=1,\ldots,p.$

    Therefore, the interpolation problem is solvable if and only if the operator $X_{\mathcal{Z},\mathcal{W}}$ admits a lift to the Schur-class interpolant.
   % Therefore the interpolation problem is equivalent to the liftability of $X_{\mathcal{Z},\mathcal{W}}$ to the interpolant. 
   % there exists $\varphi\in \mathcal{S}(\boldsymbol \theta(\Omega))$ satisfying  $\varphi(\boldsymbol \zeta_i)=w_i$ for all $i=1,\ldots,p$ if and only if the operator $X_{\mathcal{Z},\mathcal{W}}$ admits a lift to the Schur function $\varphi$.
    By Theorem \ref{1st characterization}, this occurs if and only if the functional 
    $X_{\mathcal{Q}_{\mathcal{Z}}}: \left(\mathcal{M}_{\mathcal{Q}_{\mathcal{Z}}}, \|\cdot \|_{L^1(\partial\boldsymbol \theta(\Omega))}\right) \to \mathbb C$ defined by $$X_{\mathcal{Q}_{\mathcal{Z}}}(f)=\int_{\partial\boldsymbol \theta(\Omega)} \psi f d\mu_{\boldsymbol \theta}, \qquad f\in \mathcal{M}_{\mathcal{Q}_{\mathcal{Z}}},$$
    is contractive.
    Moreover, in this contractive setting, $X_{\mathcal{Z}, \mathcal{W}} = S_{\psi} = S_{\varphi}$ on $\mathcal{Q}_{\mathcal{Z}}$, and $\psi(\boldsymbol{\zeta}_i) = w_i$ for all $i = 1, \ldots, p.$
    
    It remains only to identify $\psi.$ Since $\psi\in \mathcal{Q}_{\mathcal{Z}},$ there are scalars $c_1,\ldots,c_p$ such that $\psi=\sum_{j=1}^p c_j \mathbb S_{\boldsymbol \theta(\Omega)}(\cdot, \boldsymbol \zeta_j).$ Evaluating at each node $\boldsymbol \zeta_i$ yields
    $$ w_i=\psi(\boldsymbol \zeta_i)= \sum_{j=1}^p c_j \mathbb S_{\boldsymbol \theta(\Omega)}(\boldsymbol \zeta_i, \boldsymbol \zeta_j),\qquad i=1,\ldots,p.$$
    This uniquely determines the coefficients $c_1,\ldots, c_p$, and hence  $\psi=\psi_{\mathcal{Z,W}}.$
    \end{proof}
    
We next derive a geometric formulation of the interpolation problem using the distance characterization of commutant lifting.
%obtained in Theorem~\ref{2nd characterization}. 
As shown in the proof of Theorem~\ref{interpolation 1st}, the existence of a function $\varphi \in \mathcal{S}(\boldsymbol \theta(\Omega))$ satisfying $\varphi(\boldsymbol \zeta_i)=w_i$ for all $i=1,\ldots, p$ 
 is equivalent to the fact that the operator $X_{\mathcal{Z,W}}$ defined by \eqref{eq 12} admits a lift to $\varphi.$ In particular, such a lift necessarily implies that $\|X_{\mathcal{Z,W}}\|\leq 1.$ Applying Theorem \ref{2nd characterization} to $X_{\mathcal{Z,W}}$ therefore yields the following interpolation criterion.

%  The distance formulation of commutant lifting gives a geometric version of the preceding interpolation criterion. Let
% $X_{\mathcal Z,\mathcal W}$ be the module map defined by \eqref{eq 12}, and set
% $\psi_{\mathcal Z,\mathcal W}.$ If $\psi=0$, then $w_1=\cdots=w_p=0$, and the zero function solves the interpolation problem. Thus, in the following statement, we assume $\psi\neq0$.

\begin{theorem}\label{interpolation 2nd}
     Let $\mathcal{Z}=\{\boldsymbol \zeta_1,\ldots,\boldsymbol \zeta_p\}\subseteq \boldsymbol \theta(\Omega)$ be a set of $p$ distinct points, and let $w_1,\ldots,w_p\in \mathbb D$. Set $\psi=X_{\mathcal{Z,W}} P_{\mathcal{Q}_{\mathcal{Z}}}1$ and define $$\widetilde{\mathcal M}_{\mathcal{Q}_{\mathcal{Z}}}=\left(\mathcal{Q}_{\mathcal{Z}}^{conj}\ominus\mathrm{span}\{\overline{\psi}\}\right)\dotplus \mathcal{M}_{\boldsymbol \theta(\Omega)}\dotplus H^2_0(\partial\boldsymbol \theta(\Omega)).$$ Then there exists $\varphi\in \mathcal{S}(\boldsymbol \theta(\Omega))$ such that $\varphi(\boldsymbol \zeta_i)=w_i$ for all $i=1,\ldots, p$ if and only if 
     $$\dist_{L^1(\partial\boldsymbol \theta(\Omega))}\left(\frac{\overline{\psi}}{\|\psi\|^2_{L^2(\partial\boldsymbol \theta(\Omega))}}, \; \widetilde{\mathcal M}_{\mathcal{Q}_{\mathcal{Z}}}\right)\geq 1.$$
\end{theorem}

We now pass to quotient domains of the polydisc associated with the finite complex reflection groups $G(m,n,d)$. For these domains, Theorem \ref{3rd characterization} gives a commutant lifting criterion in terms of inner functions. Applying that criterion to the finite-dimensional zero-based quotient modules determined by the interpolation nodes yields the following inner-function characterization of Schur-class interpolation.


% In view of Theorem \ref{3rd characterization}, commutant lifting on the domain $\boldsymbol \theta(\mathbb D^d)$ is closely related to the weak-$*$ density of inner functions in the Schur class $\mathcal{S}(\boldsymbol \theta(\mathbb D^d))$. By translating this property into the interpolation setting, we can reformulate the existence of a Schur-class interpolant on $\boldsymbol \theta(\mathbb D^d)$ entirely in terms of inner functions on. This yields the following characterization.
\begin{theorem}\label{interpolation 3rd}
    Let $\mathcal{Z}=\{\boldsymbol \zeta_1,\ldots,\boldsymbol \zeta_p\}\subseteq \boldsymbol \theta(\mathbb D^d)$ be a set of $p$ distinct points, and let $w_1,\ldots,w_p\in \mathbb D.$ Then there exists $\varphi \in \mathcal{S}(\boldsymbol \theta(\mathbb D^d))$ such that $\varphi(\boldsymbol \zeta_i)=w_i$ for all $i=1,\ldots, p$ if and only if there exists a sequence $\{f_n\}\subseteq \mathcal{I}(\boldsymbol \theta(\mathbb D^d))$ satisfying $$\lim_{n \to \infty} f_n(\boldsymbol \zeta_j)=w_j, \qquad j=1,\ldots, p.$$
\end{theorem}
\begin{proof}
Let $\mathcal{Q}_{\mathcal{Z}}$ be the zero-based quotient module associated with $\mathcal{Z}$, and let $X_{\mathcal{Z,W}}\in \mathcal{B}(\mathcal{Q}_{\mathcal{Z}})$ be the module map defined by \eqref{eq 12}. Set $\psi=X_{\mathcal{Z,W}}P_{\mathcal{Q}_{\mathcal{Z}}}1.$ As in the proof of Theorem \ref{interpolation 1st}, the existence of an interpolating Schur function $\varphi$ is equivalent to $X_{\mathcal{Z,W}}$ admitting a lift to $\varphi.$ 
%function $\varphi\in \mathcal{S}(\boldsymbol \theta(\mathbb D^d))$ satisfying $\varphi(\boldsymbol \zeta_i)=w_i$ for all $i=1,\ldots,p$ is equivalent to the existence of a lift of $X_{\mathcal{Z,W}}$ to $\varphi.$ 
In this scenario, $$X_{\mathcal{Z,W}}=S_{\varphi}=S_{\psi} \quad\text{and}\quad \psi(\boldsymbol \zeta_i)=w_i \;\text{for all}\; i=1,\ldots,p.$$

Assume first that such a function $\varphi$ exists. By Theorem~\ref{3rd characterization}, there exists a sequence $\{f_n\}\subseteq \mathcal{I}(\boldsymbol \theta(\mathbb D^d))$ converging to $\psi$ in the weak-$*$ topology determined by $\mathcal{M}_{\mathcal{Q}_{\mathcal{Z}}}$.
In particular, $$\int_{\boldsymbol \theta(\mathbb T^d)} f_n \overline{\mathbb S_{\boldsymbol \theta(\mathbb D^d)}(\cdot, \boldsymbol \zeta_j)} d\mu_{\boldsymbol \theta} \to \int_{\boldsymbol \theta(\mathbb T^d)} \psi \overline{\mathbb S_{\boldsymbol \theta(\mathbb D^d)}(\cdot, \boldsymbol \zeta_j)} d\mu_{\boldsymbol \theta}, \qquad j=1,\ldots, p.$$
Applying the reproducing property, this yields
$$f_n(\boldsymbol \zeta_j)\to \psi(\boldsymbol \zeta_j)=w_j, \qquad j=1,\ldots,p.$$

Conversely, suppose that there exists a sequence $\{f_n\}\subseteq \mathcal{I}(\boldsymbol \theta(\mathbb D^d))$ such that $$ \lim_{n\to\infty} f_n(\boldsymbol{\zeta}_j)=w_j, \qquad j=1,\ldots,p.$$ 
As established prior to Theorem \ref{interpolation 1st}, the module map $X_{\mathcal{Z,W}}$ coincides with $S_{\psi}$, where $\psi= X_{\mathcal{Z,W}}P_{\mathcal{Q}_{\mathcal{Z}}}1.$ Moreover, \eqref{eq 12} and the reproducing property imply that $$\psi(\boldsymbol \zeta_j)=w_j, \qquad j=1,\ldots, p.$$
Since $\mathcal Q_{\mathcal Z}$ is spanned by the kernel functions at
$\boldsymbol\zeta_1,\ldots,\boldsymbol\zeta_p$, the reproducing property shows that $f_n$ converges to $\psi$ in the weak-$*$ topology determined by
$\mathcal Q_{\mathcal Z}^{\operatorname{conj}}$. By Corollary \ref{refined characterization}, $X_{\mathcal Z,\mathcal W}$ admits a Schur-class lift.
Hence, there exists $\varphi\in\mathcal S(\boldsymbol\theta(\mathbb D^d))$ such that $X_{\mathcal{Z,W}}=S_{\varphi}.$
Using \eqref{eq 12} and the reproducing property once more, we conclude that
% Since $$\mathcal{Q}_{\mathcal{Z}}=\mathrm{span}\{\mathbb S_{\boldsymbol \theta(\mathbb D^d)}(\cdot, \boldsymbol \zeta_i): i=1,\ldots,p\},$$ the reproducing property yields $f_n \to \psi$ in the weak-$*$ topology determined by $\mathcal{Q}_{\mathcal{Z}}^{conj}.$
%   Furthermore, both $\psi$ and $f_n$ are orthogonal to $\mathcal{M}_{\boldsymbol \theta(\mathbb D^d)}\dotplus H^2_0(\boldsymbol \theta(\mathbb T^d)).$ Consequently, $f_n \to \psi$ in the weak-$*$ topology determined by $\mathcal{M}_{\mathcal{Q}_{\mathcal{Z}}}$.
%   Theorem \ref{3rd characterization} therefore implies that $X_{\mathcal{Z,W}}$ admits a lift to some function $\varphi\in \mathcal{S}(\boldsymbol \theta(\mathbb D^d)),$ 
%     that is, $X_{\mathcal{Z,W}}=S_{\varphi}.$ Finally, arguing exactly as in the proof of Theorem \ref{interpolation 1st}, this lift satisfies 
    $$\varphi(\boldsymbol \zeta_j)=w_j, \qquad j=1,\ldots,p.$$ This completes the proof.
\end{proof}


  \subsubsection*{Acknowledgment} The author expresses his sincere gratitude to his PhD supervisor, Professor Surjit Kumar, for many valuable suggestions and comments during the preparation of this draft.


\begin{thebibliography}{33}

%\bibitem{AWY2007} A. A. Abouhajar, M. C. White, and N. J. Young, 
%A Schwarz lemma for a domain related to $\mu$-synthesis, {\it 
%J. Geom. Anal.} {\bf 17} (2007), 717-750.

%\bibitem{AJ1990} J. Agler, On the representation of certain holomorphic functions defined on a polydisc, in {\it Topics in operator theory: Ernst D. Hellinger memorial volume}, Oper. Theory Adv. Appl., {\bf 48} (1990), 47-66. 
%\bibitem{ALY2019} J. Agler, J. Lykova, N. Young, \emph{Geodesics, retracts, and the norm-preserving extension property in the symmetrized bidisc}, Mem. Amer. Math. Soc. {\bf 258} (2019), no. 1242, vii+108 pp.

\bibitem{AM1999} J. Agler, J.~E. McCarthy, \emph{Nevanlinna-Pick interpolation on the bidisk.}, J. Reine Angew. Math. {\bf 506} (1999), 191–204.

%\bibitem{AY2004} J. Agler, N. Young, \emph{The hyperbolic geometry of the symmetrized bidisc}, J. Geom. Anal. {\bf 14} (2004), no. 3, 375–403.

\bibitem{AY2017} J. Agler, N. J. Young, \emph{Realization of functions on the symmetrized bidisc}, \textit{J. Math. Anal. Appl.}, {453} (2017), 227-240.

\bibitem{AJ1995} 
J. Arazy, \emph{A survey of invariant {H}ilbert spaces of analytic functions on bounded symmetric domains}, Contemporary Mathematics, {\bf 185}  (1995), 7-65.


%\bibitem{AM2002} J. Agler and J.~E. McCarthy,  \textit{Pick interpolation and Hilbert function spaces}, Grad. Stud. Math., 44, Amer. Math. Soc., Providence, RI, 2002; MR1882259

%\bibitem{AMc2003} J. Agler and J.~E. McCarthy, Norm preserving extensions of holomorphic functions from subvarieties of the bidisk, \textit{Ann. of Math. (2)}, {\bf 157}(2003), no. 1, 289-312.

%\bibitem{AM2020}  J. Agler and J.~E. McCarthy, \textit{Operator analysis—Hilbert space methods in complex analysis}, Cambridge Tracts in Math., ({\bf 219}), Cambridge University Press, Cambridge, 2020.

%\bibitem{AY1999} J. Agler and N. J. Young, A commuting lifting theorem for a domain in $\mathbb C^2$ and spectral interpolation,  \textit{J. Funct. Anal.}, {\bf 161}(1999) 452-477.

%\bibitem{AY2017} J. Agler and N. J. Young, Realization of functions on the symmetrized bidisc, \textit{J. Math. Anal. Appl.}, {453} (2017), 227-240.

%\bibitem{Am2003} E. Amar, On the Toeplitz corona problem, \textit{Publ. Mat.}, {\bf 47}(2003), no. 2, 489-496.

%\bibitem{AT2003} C. Ambrozie, D. Timotin, A von Neumann type inequality for certain domains in $\mathbb C^n$, \textit{Proc. Amer. Math. Soc.}, {\bf 131}(2003), no. 3, 859-869.

%\bibitem{AP1990} A. V. Arkhangel’skii, L. S. Pontryagin, \textit{General topology-I}, Springer, Berlin, 1990.

%\bibitem{AV1975} W. Arveson, Interpolation problems in nest algebras, \textit{J. Funct. Anal.}, {\bf 20}(1975), no. 3, 208-233.

\bibitem{BB2004} J. A. Ball, V. Bolotnikov, \emph{Realization and interpolation for Schur-Agler-class functions on domains with matrix polynomial defining function in $\mathbb C^n$}, \textit{J. Funct. Anal.}, {\bf 213}(2004), no. 1, 45-87.

\bibitem{BLTT1999} J. A. Ball, W. S. Li, D. Timotin, T. T. Trent, 
\emph{A commutant lifting theorem on the polydisc},
Indiana Univ. Math. J. {\bf 48} (1999), no. 2, 653–675.


%\bibitem{BMV2018} J. A. Ball, G. Marx, V. Vinnikov, \emph{Interpolation and transfer-function realization for the noncommutative Schur-Agler class}, in {\it Operator theory in different settings and related applications}, Oper. Theory Adv. Appl., {\bf 262}(2018), 23-116.

\bibitem{BT1998} J. A. Ball, T. T. Trent, \emph{Unitary colligations, reproducing kernel Hilbert spaces, and Nevanlinna-Pick interpolation in several variables}, \textit{J. Funct. Anal.}, {\bf 157}(1998), no. 1, 1-61.


\bibitem{BS1982} S. R. Bell, \emph{The Bergman kernel function and proper holomorphic mappings}, Trans. Amer. Math. Soc. {\bf 270} (1982), no. 2, 685–691.

\bibitem{BDS2025} J. Bhattacharjee, K. D. Deepak, J. Sarkar, \emph{Commutant lifting and interpolation on the unit ball},  https://doi.org/10.48550/arXiv.2512.11349

\bibitem{BBS2024} T. Bhattacharyya, M. Bhowmik, H. Sau, \emph{Hankel operators and Projective Hilbert modules on quotients of bounded symmetric domains}, https://doi.org/10.48550/arXiv.2409.04582

\bibitem{TirDS2021} T. Bhattacharyya, B. Krishna Das, H. Sau, \emph{Toeplitz operators on the symmetrized bidisc}, Int. Math. Res. Not. IMRN (2021), no. 11, 8492–8520.



\bibitem{BS2018} T. Bhattacharyya, H. Sau, \emph{Holomorphic functions on the symmetrized bidisk—realization, interpolation and extension}, \textit{J. Funct. Anal.}, {\bf 274}(2018), no. 2, 504-524.


\bibitem{BDGR2022}
 S. Biswas, S. Datta, G. Ghosh, S. Shyam Roy, \emph{Reducing submodules of Hilbert modules and Chevalley-Shephard-Todd theorem}, Adv. Math. {\bf 403}(2022), 108366.

%\bibitem{BR2014} S. Biswas, S. Shyam Roy, \emph{Functional models of $\Gamma_n$- contractions and characterization of $\Gamma_n$-isometries}, J. Funct. Anal. {\bf 266} (2014), no. 10, 6224–6255.







%\bibitem{BT2014}  T. Bhattacharyya, The Tetrablock as a spectral set, \textit{Indiana Univ. Math. J.}, {\bf 63}(2014), no. 6, 1601–1629.

%\bibitem{BS2018} T. Bhattacharyya and H. Sau, Holomorphic functions on the symmetrized bidisk—realization, interpolation and extension, \textit{J. Funct. Anal.}, {\bf 274}(2018), no. 2, 504-524.

%\bibitem{BS2022} T. Bhattacharyya and H. Sau, Interpolating sequences and the Toeplitz-Corona theorem on the symmetrized bidisk, \textit{J. Operator Theory}, {\bf 87}(2022), no. 2, 435-459.

%\bibitem{BCJ2024} S. Bera, S. Chavan, S. Jain, A transference principle for involution-invariant functional Hilbert spaces, to appear in \textit{Journal of Geometric Analysis}, arXiv:2408.04384v2.

%\bibitem{CL1962} L. Carleson, Interpolations by bounded analytic functions and the corona problem, \textit{Ann. of Math. (2)}, {\bf 76}(1962), 547-559

%\bibitem{CJ2000} J. B. Conway, \textit{A course in operator theory}, Grad. Stud. Math.,{21}(2000), American Mathematical Society.


\bibitem{DS2021} B. Krishna Das, H. Sau, \emph{Algebraic properties of Toeplitz operators on the symmetrized polydisk}, Complex Anal. Oper. Theory {\bf 15} (2021), no. 3, Paper No. 60, 28 pp.

\bibitem{DeS2026} K. D. Deepak, J. Sarkar, \emph{Commuting lifting, interpolation, and perturbations on the polydisc}, to appear in Amer. J. Math.

%\bibitem{EZ2005} A. Edigarian, W. Zwonek, \emph{Geometry of the symmetrized polydisc}, Arch. Math. (Basel) {\bf 84} (2005), no. 4, 364–374.

%\bibitem{DMM2007} M.~A. Dritschel, S.~A.~M. Marcantognini and S.~A. McCullough, Interpolation in semigroupoid algebras, \textit{J. Reine Angew. Math.} {\bf 606} (2007), 1--40.

%\bibitem{DM2007} M.~A. Dritschel and S.~A. McCullough, Test functions, kernels, realizations and interpolation, in \textit{Operator theory, structured matrices, and dilations}, Theta Ser. Adv. Math., {\bf 7}(2007), 153-179.

\bibitem{GN2023} G. Ghosh, E. K. Narayanan, \emph{Toeplitz operators on the weighted Bergman spaces of quotient domains}, Bull. Sci. Math. {\bf 188} (2023), Paper No. 103340, 29 pp.

\bibitem{GR2024} G. Ghosh, S Shyam Roy,  \emph{Brown-Halmos type theorems on the proper images of bounded symmetric domains}, 
https://doi.org/10.48550/arXiv.2405.08002

\bibitem{HM1969}  K. T. Hahn, J. Mitchell, \emph{$H^p$ spaces on bounded symmetric domains},  Trans. Amer. Math. Soc. {\bf 146} (1969), 521-531. 

\bibitem{HM1973}  K. T. Hahn, J. Mitchell, \emph{$H^p$ spaces on bounded symmetric domains}, Ann. Polon. Math. {\bf 28} (1973), 89–95.

\bibitem{KA1965}
A. Koran\'yi, \emph{The Poisson integral for generalized half-planes and bounded symmetric domains}, Ann. of Math. (2),   {\bf 82} (1965), 332–350.


\bibitem{KZ2013} L. Kosi\'nski, W. Zwonek, \emph{Proper holomorphic mappings vs. peak points and Shilov boundary}, Ann. Polon. Math. {\bf 107} (2013), pp. 97–108.

\bibitem{LO1977}
O. Loos, \emph{Bounded symmetric domains and Jordan pairs}, University of California, Irvine, 1977.

%\bibitem{GS2016} S. Gorai, J. Sarkar, \emph{Characterizations of symmetrized polydisc}, Indian J. Pure Appl. Math. {\bf 47} (2016), no. 3, 391–397.

\bibitem{MRZ2013} G. Misra, S. S. Roy, G. Zhang, \emph{Reproducing kernel for a class of weighted Bergman spaces on the symmetrized polydisc}, {\it Proc. Amer. Math. Soc.} {\bf 141} (2013), no. 7, 2361-2370.

%\bibitem{NF1970} B. Sz.-Nagy and C. Foias, \emph{Harmonic analysis on operators on Hilbert spac}, North-Holland, Amsterdam, 1970.

%\bibitem{Nev1919} R. Nevanlinna, \emph{Über beschränkte Funktionen, die in gegebenen Punkten vorgeschriebene Werte annehmen}, Ann. Acad. Sci. Fenn. {\bf 13} (1919), 27-43.

%\bibitem{NN2006} N. Nikolov, \emph{The symmetrized polydisc cannot be exhausted by domains biholomorphic to convex domains}, Ann. Polon. Math. {\bf 88} (2006), no. 3, 279–283.

%\bibitem{NPTZ2008} N. Nikolov, P. Pflug, P. J. Thomas, W. Zwonek, \emph{Estimates of the Carathéodory metric on the symmetrized polydisc}, J. Math. Anal. Appl. {\bf 341} (2008), no. 1, 140–148.

%\bibitem{NPZ2007} N. Nikolov, P. Pflug, W. Zwonek, \emph{The Lempert function of the symmetrized polydisc in higher dimensions is not a distance}, Proc. Amer. Math. Soc. {\bf 135} (2007), no. 9, 2921–2928.

%\bibitem{NTT2016} N. Nikolov,  P. J. Thomas, D. Tran, \emph{Lifting maps from the symmetrized polydisc in small dimensions}, Complex Anal. Oper. Theory {\bf 10} (2016), no. 5, 921–941.

%\bibitem{Pk1916} G. Pick, \emph{Über die Beschränkungen analytischer Funktionen, welche durch vorgegebenen Funktionswerte bewiskt sind}, Math. Ann. {\bf 77} (1916), 7-23. 

\bibitem{Rd1969} W. Rudin, \emph{Function Theory in polydiscs}, W. A. Benjamin, Inc. New York-Amsterdam, 1969.

\bibitem{RD1982} W. Rudin, \emph{Proper holomorphic maps and finite reflection groups}, Indiana Univ. Math. J. {\bf 31} (1982), no. 5, 701–720.

\bibitem{Rd1983} W. Rudin, \emph{Inner functions in the unit ball of $\mathbb C^n$}, J. Funct. Anal. {\bf 50} (1983), no. 1, 100-126.

\bibitem{DS1967} D. Sarason, \emph{Generalized interpolation in $H^\infty$}, Trans. Amer. Math. soc. {\bf 127}(1967), 179-203.

\bibitem{ST1954} G. C. Shephard, J. A. Todd, \emph{Finite unitary reflection groups}, Canad. J. Math. {\bf 6} (1954), 274–304.

\bibitem{TM2013}
M. Trybula, \emph{Proper holomorphic mappings, Bell’s formula, and the Lu Qi-Keng problem on the
tetrablock}, Arch. Math. (Basel), 101 (2013), p. 549–558. 


%\bibitem{TTZ2011} P. J. Thomas, N. V. Trao, W. Zwonek, \emph{Green functions of the spectral ball and symmetrized polydisk}, J. Math. Anal. Appl. {\bf 377} (2011), no. 2, 624–630.




%\bibitem{HM1969}  K. T. Hahn, J. Mitchell, $H^p$ spaces on bounded symmetric domains, {\it Trans. Amer. Math. Soc.} {\bf 146} (1969), 521-531.



%\bibitem{PR2016}  V. Paulsen and M. Raghupathi, {\it An Introduction to the Theory of Reproducing Kernel Hilbert Spaces}, Cambridge Stud. Adv. Math., 152, Cambridge University Press, Cambridge, 2016.

%\bibitem{RD1991} W. Rudin, \textit{Functional analysis}, Internat. Ser. Pure Appl. Math., McGraw-Hill, Inc., New York, 1991.

%\bibitem{SB1978} C. F. Schubert, The corona theorem as an operator theorem, \textit{Proc. Amer. Math. Soc.}, {\bf 69}(1978), no. 1, 73-76.

%\bibitem{SZ1974} Z. Slodkowski and W. \.Zelazko, On joint spectra of commuting families of operators, \textit{Studia Math.} {\bf 50} (1974), 127-148.

%\bibitem{VS2005} V. Scheidemann, \textit{Introduction to complex analysis in several variables},  Birkhäuser Verlag, Basel, 2005.

%\bibitem{VF1982} F.-H. Vasilescu, Analytic functional calculus and spectral decompositions, {\bf 1}(1982), \textit{Math. Appl.} (East European Ser.), D. Reidel Publishing Co., Dordrecht.




% \bibitem{KW1965}
% A. Koran\'yni and J. A. Wolf, \emph{Realization of hermitian symmetric spaces as generalized half-planes}, Ann. of Math. (2)   {\bf 81} (1965), 265–288.



% \bibitem{SM1976}
% M. Stoll, \emph{The space $N_{*}$ of holomorphic functions on bounded symmetric domains}, Ann. Polon. Math. {\bf 32} (1976), no. 1, 95–110.

\bibitem{UH1996}
H. Upmeier, \emph{Toeplitz operators and index theory in several complex variables}, Operator Theory Advances and Applications,  {\bf 81}, Birkh\"auser, 1996.



\bibitem{WN1968}
N. J. Weiss, \emph{Almost everywhere convergence of Poisson integrals on generalized half-planes}, Bull. Amer. Math. Soc.,   {\bf 74} (1968), 533–537.


% \bibitem{JP2000}
% M. Jarnicki, P. Pflug, \emph{Extension of holomorphic functions}, De Gruyter Expositions in Mathematics,
% 34. De Gruyter, Berlin, 2000, x+487 pp.


% \bibitem{CC1955}
% C. Chevalley, \emph{Invariants of finite groups generated by reflections}, Amer. J. Math. {\bf 77} (1955), 778–782.




% \bibitem{EKZ2013} A. Edigarian, L. Kosi\'nski, W. Zwonek, \emph{The Lempert theorem and the tetrablock}, J. Geom. Anal. {\bf 23} (2013), no. 4, 1818–1831.


% \bibitem{BPR2012} T. Bhattacharyya, S. Pal and S. Shyam Roy, \emph{Dilations of $\Gamma$-contractions by solving operator equations}, Adv. Math. {\bf 230} (2012) 577–606.




% \bibitem{RS1977} R. P. Stanley, \emph{Relative invariants of finite groups generated by pseudoreflections}, J. Algebra, {\bf 49} (1977), p. 134–148.


% \bibitem{RD1980} W. Rudin, \emph{Function theory in the unit ball of $\mathbb C^n$}, vol. 241 of Grundlehren der Mathematischen Wissenschaften, Springer-Verlag, New York-Berlin, 1980.

%\bibitem{MRZ2013} G. Misra, S. Shyam Roy, and G. Zhang, \emph{Reproducing kernel for a class of weighted Bergman spaces on the symmetrized polydisc}, Proc. Amer. Math. Soc. {\bf 141} (2013), no. 7, 2361–2370.

%\bibitem{TirS2022} T. Bhattacharyya and H. Sau, \emph{Interpolating sequences and the Toeplitz-Corona theorem on the symmetrized bidisk}, J. Operator Theory {\bf 87} (2022), no. 2, 435–459.


\end{thebibliography}
\end{document}
